\documentclass[a4paper,11pt]{article}
\usepackage[utf8]{inputenc}
\usepackage[T2A]{fontenc}
\usepackage[english,russian]{babel}
\usepackage[margin=2.5cm]{geometry}
\usepackage{amsmath,amssymb,amsthm}
\usepackage{stmaryrd}
\usepackage{proof}

\newtheorem{thm}{Теорема}[section]
\newtheorem{lmm}[thm]{Лемма}
\newtheorem{axm}[thm]{Аксиома}
\theoremstyle{definition}
\newtheorem{dfn}[thm]{Определение}

\newcommand{\sem}[2]{\llbracket #1 \rrbracket^{#2}}
\newcommand{\cmpl}{\mathsf{c}}
\newcommand{\proves}{\vdash}
\newcommand{\True}{\text{И}}
\newcommand{\False}{\text{Л}}
\newcommand{\amp}{\mathbin{\&}}
\newcommand{\vect}[1]{\overrightarrow{#1}}
\newcommand{\gn}[1]{\text{«}#1\text{»}}
\newcommand{\Sh}[1]{\mathrm{Sh}\langle #1\rangle}
\newcommand{\Rr}[1]{\mathrm{R}\langle #1\rangle}
\newcommand{\Mn}[1]{\mu\langle #1\rangle}

\begin{document}

\section{Общие замечания}
Данный текст представляет из себя краткий конспект лекций по курсу
«Математическая логика», рассказанных студентам ИТМО (группы 2537--2539) в
2013--2014 учебном году.

\section{Исчисление высказываний}
Матлогика~--- это наука о правильных математических рассуждениях, а поскольку
рассуждения обычно ведутся на каком-то языке, то она неразрывна связана с идеей
двух языков: языка исследователя (или иначе его называют мета-языком), и
предметного языка. Как следует из названий, языком исследователя пользуемся мы,
формулируя утверждения или доказывая теоремы о разных способах математических
рассуждений, или просто их обсуждая. Сами же математические рассуждения,
собственно и составляющие предмет исследования, формализованы в некотором
предметном языке.

Мы начнём с очень простого предметного языка~--- языка исчисления высказываний.
Элементами (строками) данного языка являются некоторые выражения (формулы), по
структуре очень похожие на арифметические, которые называются высказываниями.

Каждое высказывание~--- это либо пропозициональная переменная~--- большая буква
латинского алфавита, возможно, с цифровым индексом, либо оно составлено из
одного или двух высказываний меньшего размера, соединённых логической связкой.

Связок в языке мы определим 4 (хотя при необходимости этот список может быть в
любой момент изменен).
\begin{itemize}
\item конъюнкция: если $\alpha$ и $\beta$~--- высказывания, то $\alpha \amp \beta$~---
  тоже высказывание.
\item дизъюнкция: если $\alpha$ и $\beta$~--- высказывания, то $\alpha \vee \beta$~---
  тоже высказывание.
\item импликация: если $\alpha$ и $\beta$~--- высказывания, то $\alpha \to \beta$~---
  тоже высказывание.
\item отрицание: если $\alpha$~--- высказывание, то $\neg\alpha$~--- тоже
  высказывание.
\end{itemize}
Высказывания, подробности которых нас не интересуют, мы будем обозначать
начальными буквами греческого алфавита ($\alpha$, $\beta$, $\gamma$ и т.п.).

\subsection{Оценка высказываний}
Процесс «вычисления» значения высказываний (оценка высказываний) имеет
совершенно естественное определение. Мы фиксируем некоторое множество
истинностных значений $V$, для начала мы в качестве такого множества возьмем
множество $\{\True, \False\}$, здесь $\True$ означает истину, а $\False$~---
ложь. Всем пропозициональным переменным мы приписываем некоторое значение, а
затем рекурсивно вычисляем значение выражения естественным для указанных связок
образом.

В дальнейшем мы будем брать необычные множества истинностных значений, и будем
давать неожиданный смысл связкам, однако, классическая интерпретация связок
всегда будет подразумеваться, если не указано иного.

Среди высказываний выделяются те, что остаются истинными при любой оценке
пропозициональных переменных. Такие высказывания называют тавталогиями или
общезначимыми высказываниями. Также, на языке исследователя общезначимость
высказывания $\alpha$ можно кратко записать как $\models \alpha$.

Оценку высказываний мы будем записывать с помощью двойных квадратных скобок.
Например, нетрудно видеть, что $\llbracket P \to P\rrbracket = \True$. Если нам
требуется явно задать значения некоторых пропозициональных переменных, мы будем
записывать эти значения как верхний индекс:
$\sem{P \to Q}{P:=\False} = \True$.

\subsection{Доказательства}
В любой теории есть некоторые утверждения (аксиомы), которые принимаются без
доказательства. В исчислении высказываний мы должны явно определить список всех
возможных аксиом. Например, мы можем взять утверждение $A \amp B \to A$ в
качестве аксиомы. Однако, есть множество аналогичных утверждений, например,
$B \amp A \to B$, которые не отличаясь по сути, отличаются по записи, и
формально говоря, являются другими утверждениями.

Для решения вопроса мы введём понятие схемы аксиом~--- некоторого обобщённого
шаблона, подставляя значения в который, мы получаем различные, но схожие
аксиомы. Например, схема аксиом $\psi \amp \varphi \to \psi$ позволяет получить
как аксиому $A \amp B \to A$ (при подстановке $\psi := A$, $\varphi := B$), так и
аксиому $B \amp A \to B$.

Возьмем следующие схемы аксиом для исчисления высказываний.
\begin{align*}
&(1)\ (\varphi) \to ((\psi) \to (\varphi))\\
&(2)\ ((\varphi) \to (\psi)) \to ((\varphi) \to (\psi) \to (\pi)) \to ((\varphi) \to (\pi))\\
&(3)\ (\varphi) \to (\psi) \to (\varphi) \amp (\psi)\\
&(4)\ (\varphi) \amp (\psi) \to (\varphi)\\
&(5)\ (\varphi) \amp (\psi) \to (\psi)\\
&(6)\ (\varphi) \to (\varphi) \vee (\psi)\\
&(7)\ (\psi) \to (\varphi) \vee (\psi)\\
&(8)\ ((\varphi) \to (\pi)) \to ((\psi) \to (\pi)) \to ((\varphi) \vee (\psi) \to (\pi))\\
&(9)\ ((\varphi) \to (\psi)) \to ((\varphi) \to \neg(\psi)) \to \neg(\varphi)\\
&(10)\ \neg\neg(\varphi) \to (\varphi)
\end{align*}

Помимо аксиом, нам требуется каким-то образом научиться преобразовывать одни
верные утверждения в другие. Сделаем это с помощью правил вывода. У нас пока
будет одно правило вывода~--- Modus Ponens. Это также схема, она позволяет при
доказанности двух формул $\psi$ и $\psi \to \varphi$ считать доказанной формулу
$\varphi$.

\begin{dfn}
Доказательство в исчислении высказываний~--- это некоторая конечная
последовательность выражений $\alpha_1, \alpha_2 \dots \alpha_n$ из языка $L$,
такая, что каждое из утверждений $\alpha_i$ ($1 \le i \le n$) либо является
аксиомой, либо получается из других утверждений
$\alpha_{P_1}, \alpha_{P_2} \dots \alpha_{P_k}$ ($P_1 \dots P_k < i$) по
правилу вывода.
\end{dfn}

\begin{dfn}
Высказывание $\alpha$ называется доказуемым, если существует доказательство
$\alpha_1, \alpha_2 \dots \alpha_k$, и в нем $\alpha_k$ совпадает с $\alpha$.
\end{dfn}

Вообще, схемы аксиом и правила вывода существуют для удобства задания
исчисления. В дальнейшем будет очень неудобно возиться с этими объектами.
Поэтому мы считаем, что в исчислении имеется бесконечно много аксиом и правил
вывода, которые порождаются подстановкой всех возможных формул вместо параметров
в схемы.

В качестве сокращения записи в языке исследователя мы будем писать
$\proves \alpha$, чтобы сказать, что $\alpha$ доказуемо.

Традиционно правило вывода Modus Ponens записывают так:
\[ \infer{\psi}{\varphi & (\varphi) \to (\psi)} \]

\section{Теорема о дедукции}
Соглашение об обозначениях. Будем обозначать буквами $\Gamma$, $\Delta$,
$\Sigma$, $\Pi$ списки формул (возможно, пустые).

\begin{dfn}[Вывод из допущений]
Пусть $\Gamma$~--- некоторый список высказываний, а $\alpha$~--- некоторое
высказывание. Тогда мы будем говорить, что высказывание $\alpha$ выводимо из
$\Gamma$ (и записывать это как $\Gamma \proves \alpha$), если существует такая
последовательность высказываний $\alpha_1, \alpha_2, \dots \alpha_{n-1},
\alpha$ (называемая выводом $\alpha$ из $\Gamma$), что каждое из высказываний
$\alpha_i$~--- это
\begin{itemize}
\item либо аксиома,
\item либо получается по правилу Modus Ponens из предыдущих высказываний,
\item либо~--- высказывание из списка $\Gamma$.
\end{itemize}
Элементы $\Gamma$ мы будем называть допущениями. Также эти элементы называют
предположениями или гипотезами.
\end{dfn}

В свете данного определения можно заметить, что доказательство~--- это вывод из
пустого списка допущений.

\begin{thm}[Теорема о дедукции]\label{thm:deduction}
Утверждение $\Gamma \proves \alpha \to \beta$ справедливо тогда и только тогда,
когда справедливо, что $\Gamma, \alpha \proves \beta$.
\end{thm}

Для доказательства рассмотрим следующую лемму:

\begin{lmm}
$\proves \alpha \to \alpha$.
\end{lmm}
\begin{proof}
\begin{align*}
&(1)\ \alpha \to (\alpha \to \alpha) &&\text{Сх. акс. 1}\\
&(2)\ (\alpha \to (\alpha \to \alpha)) \to (\alpha \to ((\alpha \to \alpha) \to \alpha)) \to (\alpha \to \alpha) &&\text{Сх. акс. 2}\\
&(3)\ (\alpha \to ((\alpha \to \alpha) \to \alpha)) \to (\alpha \to \alpha) &&\text{M.P. 1,2}\\
&(4)\ (\alpha \to ((\alpha \to \alpha) \to \alpha)) &&\text{Сх. акс. 1}\\
&(5)\ \alpha \to \alpha &&\text{M.P. 4,3}
\end{align*}
\end{proof}

\begin{proof}[Доказательство теоремы~\ref{thm:deduction}]
Сперва докажем прямое следствие. Для этого нам достаточно научиться по любому
выводу $\alpha \to \beta$ из $\Gamma$ строить вывод $\beta$ из $\Gamma, \alpha$.
Возьмем вывод формулы $\alpha \to \beta$, то есть некоторую последовательность
формул $\delta_1 \dots \delta_{m-1}; \alpha \to \beta$. Добавив к выводу 2
формулы, получаем требуемый вывод:
\begin{align*}
&(1)\ \delta_1\\
&\quad\dots\\
&(m-1)\ \delta_{m-1}\\
&(m)\ \alpha \to \beta\\
&(m+1)\ \alpha &&\text{«Свежедобавленная» аксиома}\\
&(m+2)\ \beta &&\text{M.P. } m, m+1
\end{align*}

Теперь докажем обратное. Нам необходимо построить вывод утверждения
$\Gamma \proves \alpha \to \beta$ по имеющемуся выводу
$\delta_1 \dots \delta_{m-1}, \beta$. Мы поступим так: сперва набросаем план
вывода~--- разместим по тексту «ключевые» формулы, которые потом дополним до
полноценного вывода промежуточными формулами.

План вывода будет такой:
\begin{align*}
&(1)\ \Gamma \proves \alpha \to \delta_1\\
&\quad\dots\\
&(m-1)\ \Gamma \proves \alpha \to \delta_{m-1}\\
&(m)\ \Gamma \proves \alpha \to \beta
\end{align*}

Теперь надо дополнить его до полноценного вывода. Будем рассматривать формулы
подряд и перед каждой формулой добавлять некоторое количество формул,
обосновывающих соответствующий шаг доказательства. Рассмотрим формулу номер $i$.
Возможны следующие варианты:
\begin{enumerate}
\item $\delta_i$~--- это аксиома или предположение, входящее в $\Gamma$. Тогда
  перед этой формулой вставим формулы $\delta_i$ и
  $\delta_i \to (\alpha \to \delta_i)$, и окажется, что $i$-я формула выводится
  из предыдущих двух формул путем применения правила Modus Ponens.
\item $\delta_i$ совпадает с $\alpha$. Тогда мы вставим перед ней 4 первые
  формулы из леммы, и $\delta_i \to \alpha$ будет получаться по правилу Modus
  Ponens.
\item $\delta_i$ выводится по правилу Modus Ponens из каких-то других
  утверждений $\delta_j$ и $\delta_k$ (при этом $\delta_k \equiv \delta_j \to
  \delta_i$), где $j < i$ и $k < i$. Покажем, что $\alpha \to \delta_i$ тоже
  может быть выведена из утверждений $\alpha \to \delta_j$ и $\alpha \to
  (\delta_j \to \delta_i)$.

  Для этого добавим два высказывания:
  \begin{align*}
  &(i-6)\ (\alpha \to \delta_j) \to ((\alpha \to (\delta_j \to \delta_i)) \to (\alpha \to \delta_i)) &&\text{Сх. акс. 2}\\
  &(i-4)\ ((\alpha \to (\delta_j \to \delta_i)) \to (\alpha \to \delta_i)) &&\text{M.P. из } j \text{ и } i-6
  \end{align*}
\end{enumerate}
\end{proof}

По аналогии мы можем рассмотреть отношение следования. Будем говорить, что
высказывание $\alpha$ следует из высказываний $\Gamma$, если при любой оценке
пропозициональных переменных, входящих в высказывания $\Gamma$ и $\alpha$, на
которых все высказывания из $\Gamma$ истинны, $\alpha$ также истинно.
Записывать, что $\alpha$ следует из $\Gamma$, будем так: $\Gamma \models \alpha$.

\section{Теорема о полноте исчисления высказываний}
\begin{dfn}
Введем обозначение. Пусть $\alpha$~--- это некоторое высказывание, а $x$~---
некоторое истинностное значение. Тогда обозначим за $[x]\,\alpha$ высказывание
$\alpha$, если $x$~--- истина, и $\neg(\alpha)$, если $x$~--- ложь. Также, если
формула $\alpha$~--- это формула с $n$ пропозициональными переменными $P_1
\dots P_n$, и $x_1 \dots x_n$~--- некоторые истинностные значения, то за
$\sem{\alpha}{P_1:=x_1,\dots P_n:=x_n}$ обозначим значение формулы $\alpha$ при
подстановке значений $x_1 \dots x_n$ вместо переменных $P_1 \dots P_n$.
\end{dfn}

\begin{lmm}\label{lmm:weak}
Если $\Gamma, \Sigma \proves \alpha$, то $\Gamma, \Delta, \Sigma \proves
\alpha$. Если $\Gamma, \Delta, \Sigma, \Pi \proves \alpha$, то $\Gamma, \Sigma,
\Delta, \Pi \proves \alpha$.
\end{lmm}
\begin{proof}
Упражнение.
\end{proof}

\begin{lmm}\label{lmm:concat}
Если справедливы 3 утверждения: $\Gamma \proves \gamma$, $\Delta \proves
\delta$ и $\gamma, \delta \proves \alpha$, то справедливо и $\Gamma, \Delta
\proves \alpha$.
\end{lmm}
\begin{proof}
Мы получим требуемый вывод, просто последовательно соединив все три исходных
вывода. Первые два вывода будут (очевидно) корректными при допущениях $\Gamma$ и
$\Delta$. В третьем же выводе могут использоваться высказывания $\gamma$ и
$\delta$, отсутствующие в предположениях. Но поскольку эти высказывания доказаны
в первых двух частях вывода, мы будем иметь полное право их упоминать~--- на тех
же основаниях, на которых они указаны в конце соответствующих доказательств.
\end{proof}

Возьмем некоторую связку исчисления высказываний, например конъюнкцию:
$A \amp B$. Построим для нее таблицу истинности. По каждой строчке построим
утверждение, в котором отрицания появляются там, где в таблице истинности
находится $\False$:
\[
\begin{array}{cc|c|l}
A & B & A \amp B & \text{утверждение}\\\hline
\False & \False & \False & \neg A, \neg B \proves \neg(A \amp B)\\
\False & \True & \False & \neg A, B \proves \neg(A \amp B)\\
\True & \False & \False & A, \neg B \proves \neg(A \amp B)\\
\True & \True & \True & A, B \proves A \amp B
\end{array}
\]

\begin{lmm}\label{lmm:truthtable}
Каждое из построенных по таблицам истинности утверждений доказуемо.
\end{lmm}
\begin{proof}
Упражнение.
\end{proof}

\begin{lmm}[Правило контрапозиции]\label{lmm:contra}
Каковы бы ни были формулы $\alpha$ и $\beta$, справедливо, что
$\proves (\alpha \to \beta) \to (\neg\beta \to \neg\alpha)$.
\end{lmm}
\begin{proof}
Сперва докажем, что $\alpha \to \beta, \neg\beta \proves \neg\alpha$.
\begin{align*}
&(1)\ (\alpha \to \beta) \to (\alpha \to \neg\beta) \to \neg\alpha &&\text{Сх. акс. 9}\\
&(2)\ \alpha \to \beta &&\text{Допущение}\\
&(3)\ (\alpha \to \neg\beta) \to \neg\alpha &&\text{M.P. 2,1}\\
&(4)\ \neg\beta \to (\alpha \to \neg\beta) &&\text{Сх. акс. 1}\\
&(5)\ \neg\beta &&\text{Допущение}\\
&(6)\ \alpha \to \neg\beta &&\text{M.P. 5,4}\\
&(7)\ \neg\alpha &&\text{M.P. 6,3}
\end{align*}
Тогда, применив 2 раза Теорему о дедукции, получим вывод требуемого
утверждения.
\end{proof}

\begin{lmm}[Правило исключенного третьего]\label{lmm:excluded}
Какова бы ни была формула $\alpha$, $\proves \alpha \vee \neg\alpha$.
\end{lmm}
\begin{proof}
Доказательство проведем в 3 этапа.
\begin{enumerate}
\item Для начала покажем $\proves \neg(\alpha\vee\neg\alpha) \to \neg\alpha$:
  \begin{align*}
  &(1)\ \alpha \to \alpha\vee\neg\alpha &&\text{Сх. акс. 6}\\
  &(2)\dots(n+1)\ \gamma_1, \dots \gamma_{n-1}, (\alpha \to \alpha\vee\neg\alpha) \to (\neg(\alpha\vee\neg\alpha) \to \neg\alpha) &&\text{Д-во из леммы~\ref{lmm:contra}}\\
  &(n+2)\ \neg(\alpha\vee\neg\alpha) \to \neg\alpha &&\text{M.P. } 1, n+1
  \end{align*}
\item Затем докажем $\proves \neg(\alpha\vee\neg\alpha) \to \neg\neg\alpha$:
  \begin{align*}
  &(1)\ \neg\alpha \to \alpha\vee\neg\alpha &&\text{Сх. акс. 7}\\
  &(2)\dots(k+1)\ \delta_1, \dots \delta_{k-1}, (\neg\alpha \to \alpha\vee\neg\alpha) \to (\neg(\alpha\vee\neg\alpha) \to \neg\neg\alpha) &&\text{Д-во из леммы~\ref{lmm:contra}}\\
  &(k+2)\ \neg(\alpha\vee\neg\alpha) \to \neg\neg\alpha &&\text{M.P. } 1, k+1
  \end{align*}
\item Теперь докажем все вместе:
  \begin{align*}
  &(1)\ \neg(\alpha\vee\neg\alpha) \to \neg\alpha &&\text{по пункту 1}\\
  &(2)\ \neg(\alpha\vee\neg\alpha) \to \neg\neg\alpha &&\text{по пункту 2}\\
  &(3)\ (\neg(\alpha\vee\neg\alpha) \to \neg\alpha) \to (\neg(\alpha\vee\neg\alpha) \to \neg\neg\alpha) \to (\neg\neg(\alpha\vee\neg\alpha)) &&\text{Сх. акс. 9}\\
  &(4)\ (\neg(\alpha\vee\neg\alpha) \to \neg\neg\alpha) \to \neg\neg(\alpha\vee\neg\alpha) &&\text{M.P. 1,3}\\
  &(5)\ \neg\neg(\alpha\vee\neg\alpha) &&\text{M.P. 2,4}\\
  &(6)\ \neg\neg(\alpha\vee\neg\alpha) \to (\alpha\vee\neg\alpha) &&\text{Сх. акс. 10}\\
  &(7)\ \alpha\vee\neg\alpha &&\text{M.P. 5,6}
  \end{align*}
\end{enumerate}
\end{proof}

\begin{lmm}[Об исключении допущения]\label{lmm:elim}
Пусть справедливо $\Gamma, \rho \proves \alpha$ и $\Gamma, \neg\rho \proves
\alpha$. Тогда также справедливо $\Gamma \proves \alpha$.
\end{lmm}
\begin{proof}
Применив теорему о дедукции к условиям теоремы получим следующее:
\begin{align*}
&\Gamma \proves \rho \to \alpha\\
&\Gamma \proves \neg\rho \to \alpha
\end{align*}
Тогда следующий вывод покажет $\Gamma \proves \alpha$:
\begin{align*}
&(1)\dots(p)\ \gamma_1, \dots \gamma_{p-1}, \rho \to \alpha &&\text{Вывод } \Gamma \proves \rho \to \alpha\\
&(p+1)\dots(q)\ \delta_1, \dots \delta_{q-p-1}, \neg\rho \to \alpha &&\text{Вывод } \Gamma \proves \neg\rho \to \alpha\\
&(q+1)\dots(r)\ \epsilon_1, \dots \epsilon_{r-q-1}, \rho \vee \neg\rho &&\text{Лемма~\ref{lmm:excluded}}\\
&(r+1)\ (\rho \to \alpha) \to (\neg\rho \to \alpha) \to (\rho \vee \neg\rho) \to \alpha &&\text{Сх. аксиом 8}\\
&(r+2)\ (\neg\rho \to \alpha) \to (\rho \vee \neg\rho \to \alpha) &&\text{M.P. } 1, r+1\\
&(r+3)\ \rho \vee \neg\rho \to \alpha &&\text{M.P. } p, r+2\\
&(r+4)\ \alpha &&\text{M.P. } r, r+3
\end{align*}
\end{proof}

\begin{thm}[О полноте исчисления высказываний]
Пусть справедливо $\models \alpha$. Тогда также справедливо, что $\proves
\alpha$.
\end{thm}
\begin{proof}
Для доказательства теоремы мы докажем чуть более сильное утверждение~--- что
для любого $k$ от $0$ до $n$ и любой оценки переменных $x_1, \dots x_k$
справедливо $[x_1]P_1, \dots [x_k]P_k \proves \alpha$. Нетрудно заметить, что
утверждение теоремы непосредственно следует из данного утверждения для $k = 0$.
Доказательство будет вестись индукцией по $n - k$.

\emph{База.} Пусть $n - k = 0$, то есть $k = n$. $\models \alpha$ означает, что
при любой оценке $x_1, \dots x_n$ пропозициональных переменных $P_1, \dots P_n$
справедливо $\alpha[P_1 := x_1, \dots P_n := x_n] = \True$. Возьмем некоторую
оценку переменных $x_1, \dots x_n$. Тогда, по лемме~\ref{lmm:truthtable},
$[x_1]P_1, \dots [x_n]P_n \proves [\sem{\alpha}{P_1:=x_1,\dots P_n:=x_n}]\,\alpha$,
то есть $[x_1]P_1, \dots [x_n]P_n \proves \alpha$.

\emph{Переход.} Пусть утверждение уже доказано для некоторого $n - k > 0$,
покажем его для $n - k + 1$ (то есть доказано для $k \le n$, покажем его для
$k - 1$). Возьмем некоторую оценку переменных $x_1, \dots x_{k-1}$. По
предположению, $[x_1]P_1, \dots [x_k]P_k \proves \alpha$, то есть
\begin{align*}
&[x_1]P_1, \dots [x_{k-1}]P_{k-1}, \neg P_k \proves \alpha\\
&[x_1]P_1, \dots [x_{k-1}]P_{k-1}, P_k \proves \alpha
\end{align*}
Тогда по лемме об исключении допущения, справедливо $[x_1]P_1, \dots
[x_{k-1}]P_{k-1} \proves \alpha$.
\end{proof}

\begin{thm}[О корректности исчисления высказываний]
Пусть справедливо $\proves \alpha$. Тогда также справедливо, что $\models
\alpha$.
\end{thm}
\begin{proof}
По условию теоремы, у нас есть доказательство высказывания $\alpha$, то есть
последовательность высказываний $\alpha_1, \dots \alpha_m$. Каждое
высказывание~--- это либо аксиома, либо применение правила Modus Ponens.
Докажем, что для каждого $k$ все высказывания $\alpha_l$ при $l \le k$~---
тавтологии. Доказательство будем вести индукцией по $k$.

\emph{База.} Пусть $k = 0$, тогда нет ни одного высказывания, про которое нужно
доказать, что оно~--- тавтология, то есть утверждение автоматически верно.

\emph{Переход.} Пусть для некоторого $k$ утверждение справедливо, докажем его
для $k+1$. Выберем некоторую оценку $x_1, \dots x_n$ пропозициональных
переменных $P_1, \dots P_n$, использованных в высказываниях $\alpha_1 \dots
\alpha_{k+1}$. Рассмотрим случаи.

Пусть $\alpha_{k+1}$~--- аксиома. В данную аксиому входят одна, две или три
формулы $\beta_1, \beta_2, \beta_3$. Подставив все возможные истинностные
значения вместо данных формул, можно проверить, что все аксиомы являются
тавтологиями, значит, они будут истинны и на тех конкретных значениях, которые
примут данные формулы после подстановки значений $x_1, \dots x_n$.

Пусть $\alpha_{k+1}$ получается по правилу Modus Ponens из $\alpha_p$ и
$\alpha_q$, причем $\alpha_q \equiv \alpha_p \to \alpha_{k+1}$. Тогда
$\sem{\alpha_p}{P_1:=x_1,\dots P_n:=x_n} = \True$ и $\sem{\alpha_p \to
\alpha_{k+1}}{P_1:=x_1,\dots P_n:=x_n} = \True$. Из таблицы истинности
импликации следует, что неизбежно $\sem{\alpha_{k+1}}{P_1:=x_1,\dots
P_n:=x_n} = \True$.
\end{proof}

Заметим, что вместе из этих двух теорем следует, что если неверно, что $\proves
\alpha$, то неизбежно найдется контрпример.

\section{Интуиционистское исчисление высказываний}
Заменим аксиому устранения двойного отрицания на $\varphi \to \neg\varphi \to
\psi$. Полученная система называется интуиционистским исчислением
высказываний.

\begin{thm}
В интуиционистском исчислении высказываний невозможно доказать правило
исключенного третьего: $P \vee \neg P$.
\end{thm}

Мы рассмотрим модельное доказательство. То есть, мы построим некоторую модель
для исчисления высказываний~--- так определим истинностные значения и связки,
что в нем правило исключенного третьего перестанет быть тавтологией, при этом
построенная нами модель останется корректной.

Пусть множество истинностных значений будет состоять из трех элементов: $\True$,
Н и $\False$. Определим отношение порядка, будем считать, что $\True > \text{Н}
> \False$.

Все связки на истине и лжи действуют традиционно, необходимо описать только их
поведение с элементом Н. Для двух операций дадим естественные обобщения.
$A \amp B = \min(A, B)$, $A \vee B = \max(A, B)$. Однако, положим
$\neg\text{Н} = \False$. Также, операцию $A \to B$ положим $\neg A \vee B$, за
исключением $\text{Н} \to \text{Н} = \True$.

Будем считать, что высказывание истинно на некоторой оценке, если оно на ней
принимает значение $\True$. Общезначимость же ($\models \alpha$) означает, что
$\alpha[x_1, \dots x_n] = \True$ для любой оценки переменных $x_1, \dots x_n$.

Нетрудно показать, что так определенная модель исчисления является корректной~---
то есть любое утверждение, имеющее доказательство, является тавтологией в смысле
данной модели.

\begin{proof}[Доказательство теоремы]
Поскольку $\text{Н} \vee \neg\text{Н} = \text{Н}$, то в данной модели неверно
$\models P \vee \neg P$. Значит, если бы существовало доказательство этого
утверждения, то рассмотренная модель была бы некорректной. Значит, для
интуиционистского исчисления высказываний неверно $\proves P \vee \neg P$.
\end{proof}

Однако, так построенная логика (с такой моделью) не является полной. Например,
высказывание $\neg P \vee \neg\neg P$, не являющееся доказуемым в
интуиционистской логике, тем не менее оказывается истинным. Сейчас мы рассмотрим
другие модели, обладающие свойством полноты.

\subsection{Топологическая интерпретация интуиционистского исчисления высказываний}
Пусть дано некоторое исчисление высказываний, для которого нам нужно построить
модель~--- предложить способ оценки истинности выражений. Начинаем мы с
множества истинностных значений. Возьмем в качестве этого множества все
открытые множества некоторого заранее выбранного топологического пространства.
Определим оценку для связок интуиционистского исчисления высказываний следующим
образом:
\begin{align*}
\llbracket A \amp B\rrbracket &= \llbracket A\rrbracket \cap \llbracket B\rrbracket\\
\llbracket A \vee B\rrbracket &= \llbracket A\rrbracket \cup \llbracket B\rrbracket\\
\llbracket A \to B\rrbracket &= (\cmpl\llbracket A\rrbracket \cup \llbracket B\rrbracket)^\circ\\
\llbracket \neg A\rrbracket &= (\cmpl\llbracket A\rrbracket)^\circ
\end{align*}
Будем считать, что формула истинна в данной модели, если её значение оказалось
равно всему пространству.

Например, возьмем в качестве пространства $\mathbb{R}$, и вычислим значение
формулы $A \vee \neg A$ при $A$ равном $(0, 1)$:
\begin{multline*}
\llbracket A\vee\neg A\rrbracket = (0,1)\cup\llbracket\neg A\rrbracket = (0,1)\cup(\cmpl(0,1))^\circ\\
= (0,1)\cup((-\infty,0)\cup(1,\infty)) = (-\infty,0)\cup(0,1)\cup(1,\infty).
\end{multline*}
Нетрудно видеть, что данная формула оказалась не общезначимой в данной
интерпретации.

\subsection{Модели Крипке}
См.\ Шень, Верещагин.

\section{Исчисление предикатов}
Очень мало логических утверждений истинны или ложны сами по себе. Часто
истинность утверждений зависит от каких-то обстоятельств. Для формализации этого
соображения в математической логике рассматривают понятия предметного множества
и предиката.

Сперва объясним эти понятия на языке теории моделей: если мы строим теорию о
некоторых объектах (например, о натуральных числах или о множествах), то
изучаемые объекты все вместе образуют некоторое предметное множество (будем его
обозначать буквой $D$). На данном множестве могут быть заданы функции,
отображающие элементы $D$ в $D$, и предикаты~--- отображающие элементы $D$ в
истинностные значения.

Например, если мы изучаем арифметические выражения над целыми числами, то
предметным множеством будет множество целых чисел, сложение и вычитание~--- это
двуместные функции, константу $0$ мы можем при желании рассмотреть, как
нуль-местную функцию, а отношения $(<)$ или $(=)$~--- это двуместные предикаты.

Теперь рассмотрим эту же конструкцию более формально, с точки зрения теории
доказательств.

\begin{dfn}
Терм исчисления предикатов (еще мы его будем называть предметным выражением)~---
это:
\begin{itemize}
\item предметная переменная~--- маленькая буква начала или конца латинского
  алфавита, возможно, с индексом или апострофом.
\item применение функции (функции мы будем обозначать латинскими буквами
  середины алфавита: $f, g, h, \dots$): если $\theta_1 \dots \theta_n$~--- термы
  и $f$~--- функциональный символ (то есть символ, обозначающий некоторую
  функцию), то $f(\theta_1, \dots \theta_n)$~--- тоже терм. Также, частным
  случаем функций являются константы~--- это нульместные функции. Обычно мы их
  будем обозначать маленькими буквами $c$ или $d$, возможно, с индексами.
\end{itemize}
\end{dfn}

\begin{dfn}
Формула исчисления предикатов (еще мы ее будем называть логическим выражением)~---
это:
\begin{itemize}
\item если $\alpha$ и $\beta$~--- формулы исчисления предикатов, то $\neg\alpha$,
  $\alpha \amp \beta$, $\alpha \vee \beta$, $\alpha \to \beta$~--- также
  формула. Связки расставлены в порядке убывания приоритета. Как и в исчислении
  высказываний, импликация правоассоциативна, остальные операции~---
  левоассоциативны.
\item если $\alpha$~--- формула и $x$~--- предметная переменная, то
  $\forall x\,\alpha$ и $\exists x\,\alpha$~--- также формулы. Кванторы имеют
  приоритет, одинаковый с отрицанием, и, как и отрицание, действуют только на
  ближайшее за ними логическое выражение. Например, формула
  $\forall x P(x, 5) \vee P(x, 7)$ соответствует формуле
  $(\forall x P(x, 5)) \vee P(x, 7)$.
\item применение предиката (предикаты мы будем обозначать большими латинскими
  буквами, возможно, с индексами): если $\theta_1 \dots \theta_n$~--- термы, и
  $P$~--- предикатный символ, то $P(\theta_1, \dots \theta_n)$~--- формула. В
  частности, при $n = 0$ предикат становится аналогом предметной переменной из
  исчисления высказываний.
\end{itemize}
\end{dfn}

Обратите внимание на понятия предикатный символ и функциональный символ.
Предикат и функция~--- это обычные математические функции, отображающие $D^n$ в
$V$ или $D$. Предикатные и функциональные символы же~--- это некоторые значки в
формуле. Им можно сопоставить различные предикаты и функции, в зависимости от
обстоятельств. Мы иногда будем для краткости опускать слово «символ», когда
будем говорить о тексте формул, но это различие всегда надо иметь в виду.

Также заметим, что ничего не мешает нам сделать синтаксис выражений более
приближенным к традиционному при необходимости (это будет, естественно, всякий
раз оговариваться). Например, в выражении $5 = x + 2$, понимаемом в смысле
исчисления предикатов, имеется один предикатный символ $(=)$ и три
функциональных символа: два нульместных~--- $2$, $5$~--- и один двуместный~---
$(+)$.

\subsection{Доказательства в исчислении предикатов}
\begin{dfn}
Дана некоторая формула $s$. Будем говорить, что подстрока $s_1$ строки $s$
является подформулой, если она в точности соответствует какому-то одному
нетерминалу в дереве разбора строки $s$.
\end{dfn}

\begin{dfn}
Если в формулу входит подформула, полученная по правилам для кванторов (то есть,
$\forall x\,\alpha$ или $\exists x\,\alpha$), то мы будем говорить, что формула
$\alpha$ находится в области действия данного квантора по переменной $x$. Также,
будем говорить, что любая подформула формулы $\alpha$ находится в области
действия данного квантора.
\end{dfn}

\begin{dfn}
Если некоторое вхождение переменной $x$ находится в области действия квантора по
переменной $x$, то такое вхождение мы назовем связанным. Вхождение переменной
$x$ непосредственно рядом с квантором $(\forall x \dots)$ мы назовем
связывающим. Те вхождения переменных, которые не являются связанными или
связывающими, назовем свободными. Формула, не имеющая свободных вхождений
переменных, называется замкнутой.
\end{dfn}

\begin{dfn}
Будем говорить, что терм $\theta$ свободен для подстановки в формулу $\psi$
вместо $x$ (или просто свободен для подстановки вместо $x$), если после
подстановки $\theta$ вместо свободных вхождений $x$ ни одно вхождение свободной
переменной в $\theta$ не станет связанным.
\end{dfn}

Чтобы получить список аксиом для исчисления предикатов, возьмем все схемы аксиом
исчисления высказываний и дополним их следующими двумя схемами. Здесь $x$~---
переменная, $\psi$~--- некоторая формула, $y$~--- некоторая переменная. Запись
$\psi[x := \theta]$ будет означать результат подстановки $\theta$ в $\psi$ вместо
всех свободных вхождений $x$. Пусть $\theta$ свободно для подстановки вместо
$x$.
\begin{align*}
&(11)\ \forall x(\psi) \to (\psi[x := \theta])\\
&(12)\ (\psi[x := \theta]) \to \exists x(\psi)
\end{align*}
Заметим, что если взять формулу $\exists x A(x, y)$, то по схеме аксиом (11), если
игнорировать ограничение на свободу для подстановки, следующее утверждение
должно быть тавтологией: $\forall y\exists x A(x, y) \to \exists x A(x, x)$.
Однако, оно ей не является.

Все аксиомы, порожденные данными схемами в новом языке, мы назовем аксиомами
исчисления предикатов.

\emph{Правила вывода.} Пусть $x$ не входит свободно в $\varphi$. Тогда
рассмотрим следующие дополнительные правила вывода исчисления предикатов:
\[
\infer{(\varphi) \to \forall x(\psi)}{(\varphi) \to (\psi)}
\qquad
\infer{\exists x(\psi) \to (\varphi)}{(\psi) \to (\varphi)}
\]
Добавив эти схемы к схеме для правила Modus Ponens исчисления высказываний, мы
сможем породить множество правил вывода.

\subsection{Оценка выражений в исчислении предикатов}
Для задания оценки для выражения в исчислении предикатов необходимо вместо
оценки для пропозициональных переменных $v_P$ в исчислении высказываний ввести
оценку для предикатных символов: для каждого $k$-местного предикатного символа
$P^k_n$ определить функцию $v_{P^k_n} : D^k \to V$. Аналогично, для каждого
$k$-местного функционального символа $f^k_n$ также надо определить функцию
$v_{f^k_n}$.

\begin{dfn}
Формула в исчислении предикатов общезначима, если она истинна на любом
предметном множестве $D$, при любой оценке предикатных и функциональных
символов, и при любых оценках свободных предметных переменных.
\end{dfn}

\begin{dfn}
Пусть имеется некоторое исчисление предикатов с множеством аксиом $A$, и пусть
дан некоторый (возможно, пустой) список $\Gamma$ формул исчисления предикатов.
Тогда, вывод формулы $\alpha$ в исчислении с аксиомами $A \cup \Gamma$ мы
назовем выводом из допущений $\Gamma$, и будем записывать это как $\Gamma
\proves \alpha$.
\end{dfn}

Формулы в списке формул не обязаны быть замкнутыми: например, легко показать,
что $P(x, y), P(x, y) \to P(y, x) \proves P(y, x)$. Однако, со свободными
переменными в допущениях надо быть осторожными, что находит своё отражение в
теореме о дедукции для исчисления предикатов.

\begin{thm}[Теорема о дедукции]
Если $\Gamma, \alpha \proves \beta$, и в доказательстве отстутствуют применения
правил для кванторов, использующих свободные переменные из формулы $\alpha$, то
$\Gamma \proves \alpha \to \beta$. Обратно, если $\Gamma \proves \alpha \to
\beta$, то $\Gamma, \alpha \proves \beta$.
\end{thm}
\begin{proof}
Доказательство обратного утверждения аналогично доказательству для теоремы о
дедукции для исчисления высказываний.

Доказательство прямого утверждения также проведём аналогично теореме о дедукции
для исчисления высказываний~--- индукцией и разбором случаев. 3 старых случая
остаются те же, что и в исчислении высказываний, добавилось 2 новых правила
вывода. Упражнение.
\end{proof}

\begin{lmm}\label{lmm:subst}
$\sem{\psi}{x:=\llbracket\theta\rrbracket} = \llbracket\psi[x := \theta]\rrbracket$,
если $\theta$ свободна для подстановки в $\psi$ вместо $x$.
\end{lmm}
\begin{proof}
Воспользуемся структурной индукцией по выражению $\psi$.

\emph{База индукции.} Пусть $\psi$~--- это предикат. Тогда из способа вычисления
значения предикатов очевидно, что $\llbracket P(\theta, y_1, \dots y_n)\rrbracket
= \sem{P(x, y_1, \dots y_n)}{x:=\llbracket\theta\rrbracket}$.

\emph{Переход индукции.} Пусть $\psi$ составлено из одного или двух выражений,
про которые утверждение уже доказано. Рассмотрим варианты:
\begin{itemize}
\item $\psi \equiv \alpha \to \beta$, $\psi \equiv \alpha \amp \beta$,
  $\psi \equiv \alpha \vee \beta$ или $\psi \equiv \neg\alpha$. Рассмотрение всех
  этих случаев схоже, возьмем для примера конъюнкцию.
  \begin{multline*}
  \sem{\alpha \amp \beta}{x:=\llbracket\theta\rrbracket}
  = \sem{\alpha}{x:=\llbracket\theta\rrbracket} \amp \sem{\beta}{x:=\llbracket\theta\rrbracket}
  = \llbracket\alpha[x := \theta]\rrbracket \amp \llbracket\beta[x := \theta]\rrbracket\\
  = \llbracket\alpha[x := \theta] \amp \beta[x := \theta]\rrbracket
  = \llbracket(\alpha \amp \beta)[x := \theta]\rrbracket
  \end{multline*}
\item $\psi \equiv \forall y\,\alpha$ или $\psi \equiv \exists y\,\alpha$. Опять
  же, рассмотрение случаев аналогично. Покажем, что
  $\sem{\forall y\,\alpha}{x:=\llbracket\theta\rrbracket} =
  \llbracket(\forall y\,\alpha)[x := \theta]\rrbracket$.

  Без уменьшения общности можем предположить, что $x$ входит свободно в
  $\forall y\,\alpha$~--- иначе данная замена не оказывает влияние на вычисление
  результата, и равенство очевидно.

  Поскольку $\theta$ свободна для подстановки вместо $x$ в $\psi$, то, значит,
  $y$ не входит свободно в $\theta$ (иначе свободный $y$ стал бы связанным при
  подстановке, а подстановки обязательно будут иметь место, раз $x$ входит
  свободно в $\forall y\,\alpha$).

  По предположению индукции, $\sem{\alpha}{x:=\llbracket\theta\rrbracket} =
  \llbracket\alpha[x := \theta]\rrbracket$. И, более того,
  $\sem{\alpha}{x:=\llbracket\theta\rrbracket,\,y:=v_y} =
  \sem{\alpha[x := \theta]}{y:=v_y}$, поскольку вычисление оценки $\theta$ не
  зависит от значения переменной $y$. Значит, оценки полных формул с квантором
  также будут совпадать.
\end{itemize}
\end{proof}

\begin{thm}
Исчисление предикатов корректно, т.е.\ любое доказуемое утверждение
общезначимо.
\end{thm}
\begin{proof}
Рассмотрим некоторое доказательство $\gamma_1, \dots \gamma_n$, и покажем, что
каждое из утверждений в доказательстве является общезначимым. Этим мы покажем
утверждение теоремы.

Рассмотрим какое-то утверждение $\gamma_i$. Оно либо является частным случаем
какой-либо из схем аксиом, либо получено из предыдущих утверждений
доказательства путем применения правил вывода.

Мы ограничимся рассмотрением только новых (по сравнению с исчислением
высказываний) схем аксиом и правил, поскольку только в них доказательство будет
отличаться от аналогичного для исчисления высказываний.

Возможны следующие варианты:
\begin{enumerate}
\item $\gamma_i \equiv \forall x(\psi) \to (\psi[x := \theta])$ при некоторых $x$
  (имя переменной тоже может быть разным), $\psi$ и $\theta$. Значит, нам
  необходимо показать, что при любых $x$, $\psi$ и $\theta$ (естественно, таких,
  что $\theta$ свободна для подстановки вместо $x$ в $\psi$) данное выражение
  будет истинным в любой оценке.

  Фиксируем некоторую оценку и докажем утверждение от противного. Пусть это не
  так, и $\llbracket\forall x(\psi) \to (\psi[x := \theta])\rrbracket = \False$.
  Тогда неизбежно, что $\llbracket\forall x(\psi)\rrbracket = \True$, но
  $\llbracket\psi[x := \theta]\rrbracket = \False$ (поскольку при остальных
  значениях подвыражений импликация будет истинна).

  Пусть в данной оценке $\llbracket\theta\rrbracket = v_\theta$. По
  лемме~\ref{lmm:subst} верно, что $\llbracket\psi[x := \theta]\rrbracket =
  \sem{\psi}{x:=v_\theta}$, причём, поскольку $\llbracket\psi[x :=
  \theta]\rrbracket = \False$, то $\sem{\psi}{x:=v_\theta} = \False$. Но тогда
  $\llbracket\forall x(\psi)\rrbracket = \True$ означает, что при любом $v_x \in
  D$ будет выполнено $\sem{\psi}{x:=v_x} = \True$. Противоречие.
\item $\gamma_i \equiv (\psi[x := \theta]) \to \exists x(\psi)$ при некоторых
  $x$, $\psi$ и $\theta$. Рассмотрение аналогично предыдущему пункту: заметим,
  что импликация опровергается только если $\llbracket\exists x(\psi)\rrbracket =
  \False$, но $\llbracket\psi[x := \theta]\rrbracket = \True$. По
  лемме~\ref{lmm:subst} при вычислении оценки формулы мы можем отдельно
  вычислить $v_\theta = \llbracket\theta\rrbracket$, и показать, что
  $\llbracket\psi[x := \theta]\rrbracket = \sem{\psi}{x:=v_\theta}$, что
  гарантирует существование элемента в $D$, на котором формула $\psi$ истинна.
\item $\gamma_i \equiv \varphi \to \forall x\,\psi$ и существует такой $j < i$,
  что $\gamma_j \equiv \varphi \to \psi$, при этом $x$ не входит свободно в
  $\varphi$.

  Опять же, фиксируем некоторую оценку. Нам нужно показать, что $\forall
  x\,\psi$ не бывает ложным, если $\varphi$ истинно.

  По предположению индукции мы знаем, что $\llbracket\varphi \to \psi\rrbracket =
  \True$. Пусть $\llbracket\forall x\,\psi\rrbracket = \False$. То есть, найдётся
  такой $v_x$ из $D$, что $\sem{\psi}{x:=v_x} = \False$. Но заметим, что $x$ не
  входит свободно в $\varphi$. Значит, $\sem{\varphi}{x:=v_x} = \True$. Значит,
  $\sem{\varphi \to \psi}{x:=v_x} = \False$. Противоречие с общезначимостью
  $\gamma_j$.
\item Правило для квантора существования доказывается аналогично.
\end{enumerate}
\end{proof}

\section{Полнота исчисления предикатов}
Доказательство этого факта довольно объемно, поэтому мы разделим его на три
части.
\begin{enumerate}
\item Мы научимся работать с произвольными моделями, уйдя от чрезмерного
  разнообразия возможных предметных множеств, предикатов и функциональных
  символов с помощью понятия непротиворечивого множества формул.
\item Мы покажем полноту бескванторной части исчисления предикатов.
\item Мы сведем полноту исчисления предикатов к полноте бескванторной его
  части.
\end{enumerate}

\subsection{Непротиворечивое множество формул}
Заметим, что само по себе предметное множество в анализе полноты исчисления
фигурирует в стороне от главного вопроса. После того, как мы вычислили
истинность или ложность конкретных формул в данной модели, мы о конкретном
предметном множестве забываем. Соответственно, каждая модель для нашей цели
исчерпывающе описывается списком формул, которые в этой модели истинны (в
частности, если две разных модели дают одинаковую оценку для каждой из формул~---
нам нет смысла эти модели разделять). Поэтому мы можем изучать свойства оценок и
моделей не прямо, а посредством набора истинных формул. Следующие определения
формализуют это понятие.

\begin{dfn}
Назовём $\Gamma$~--- множество замкнутых формул~--- непротиворечивым, если ни
для какой формулы $\alpha$ невозможно показать, что $\Gamma \proves \alpha$ и
$\Gamma \proves \neg\alpha$.
\end{dfn}

\begin{dfn}
Полным непротиворечивым множеством (непротиворечивым бескванторным множеством)
формул назвем такое множество $\Gamma$, что для любой замкнутой (замкнутой и
бескванторной) формулы $\alpha$ либо $\alpha \in \Gamma$, либо $(\neg\alpha) \in
\Gamma$.
\end{dfn}

\begin{lmm}\label{lmm:consistent}
Если $\Gamma$~--- непротиворечивое множество замкнутых (бескванторных) формул,
то для любой замкнутой (бексванторной) формулы $\alpha$ либо $\Gamma \cup
\{\alpha\}$, либо $\Gamma \cup \{\neg\alpha\}$ непротиворечиво.
\end{lmm}
\begin{proof}
Пусть это не так, и найдутся такие $\beta$ и $\delta$, что $\Gamma, \alpha
\proves \beta \amp \neg\beta$ и $\Gamma, \neg\alpha \proves \delta \amp
\neg\delta$. Без ограничения общности мы можем предположить, что $\beta \equiv
\delta$, поскольку если мы показали $\beta \amp \neg\beta$, то мы можем показать
и $\delta \amp \neg\delta$ (это можно показать на основании доказуемости
формулы $\psi \to \neg\psi \to \varphi$).

Тогда рассмотрим следующее доказательство в предположении $\Gamma$:
\begin{align*}
&(1\dots k)\ \alpha \to \beta \amp \neg\beta &&\text{Т. о дедукции}\\
&(k+1\dots l)\ \neg\alpha \to \beta \amp \neg\beta &&\text{Т. о дедукции}\\
&(l+1)\ (\alpha \to \beta \amp \neg\beta) \to (\neg\alpha \to \beta \amp \neg\beta) \to (\alpha \vee \neg\alpha \to \beta \amp \neg\beta) &&\text{Сх. акс. 8}\\
&(l+2)\ (\neg\alpha \to \beta \amp \neg\beta) \to (\alpha \vee \neg\alpha \to \beta \amp \neg\beta) &&\text{M.P. } k, l+1\\
&(l+3)\ \alpha \vee \neg\alpha \to \beta \amp \neg\beta &&\text{M.P. } l, l+2\\
&(l+4\dots m)\ \alpha \vee \neg\alpha &&\text{Лемма~\ref{lmm:excluded}}\\
&(m+1)\ \beta \amp \neg\beta &&\text{M.P. } m, l+3
\end{align*}
Таким образом, имея доказательства противоречивости $\Gamma, \alpha$ и $\Gamma,
\neg\alpha$, мы можем построить доказательство противоречивости и самого
$\Gamma$.
\end{proof}

\begin{thm}\label{thm:complete-set}
Любое множество непротиворечивых формул $\Gamma$ мы можем дополнить до полного
(полного бескванторного) множества.
\end{thm}
\begin{proof}
Упорядочим все возможные формулы (бескванторные формулы) исчисления (их, как не
трудно заметить, счётное количество, и мы можем их занумеровать целыми числами):
$\gamma_1, \gamma_2, \dots$. По данной последовательности построим
последовательность множеств $\Gamma_1, \Gamma_2, \dots$. Положим $\Gamma_1 =
\Gamma$. Рассмотрим некоторую формулу $\gamma_n$. По предыдущей лемме, либо
$\Gamma \cup \{\gamma_n\}$, либо $\Gamma \cup \{\neg\gamma_n\}$ непротиворечиво.
Пусть для определенности это $\Gamma \cup \{\gamma_n\}$. Тогда положим
$\Gamma_{n+1} = \Gamma_n \cup \{\gamma_n\}$.

Возьмем множество $\Gamma^* = \bigcup \Gamma_n$. Ясно, что это множество
полное~--- поскольку мы перебрали все формулы и рассматривали каждую формулу
вместе со своим отрицанием. Также ясно, что оно непротиворечиво: иначе есть
доказательство противоречия (оно, естественно, конечного размера),
использующего формулы $\gamma_{p_1} \dots \gamma_{p_n}$, каждая из которых
добавлена на каком-то шаге. Но тогда и множество $\Gamma_{\max(p_1,\dots
p_n)+1}$~--- множество, построенное при добавлении последней из формул
$\gamma_{p_i}$~--- противоречиво, что доказывает утверждение.
\end{proof}

\begin{dfn}
Моделью непротиворечивого множества формул мы назовем такие оценки предикатов и
функциональных символов, что каждая из формул данного множества истинна. Также,
по аналогии с исчислением высказываний, введём обозначение: $\Gamma \models
\alpha$ ($\alpha$ следует из $\Gamma$), если $\llbracket\alpha\rrbracket =
\True$ в любой модели множества $\Gamma$.
\end{dfn}

\begin{thm}
Если $\Gamma \proves \alpha$, то $\Gamma \models \alpha$.
\end{thm}
\begin{proof}
Механическая проверка всех правил и схем аксиом.
\end{proof}

\begin{thm}
Если $\Gamma$ имеет модель, то оно непротиворечиво.
\end{thm}
\begin{proof}
Пусть это не так, то есть $\Gamma \proves \alpha$ и $\Gamma \proves \neg\alpha$.
Значит, $\Gamma \models \alpha$ и $\Gamma \models \neg\alpha$. То есть,
$\llbracket\alpha\rrbracket = \True$. Значит, по определению оценки для
отрицания, $\llbracket\neg\alpha\rrbracket = \False$. Но это вступает в
противоречие с $\Gamma \models \neg\alpha$.
\end{proof}

\subsection{Полнота бескванторной части исчисления предикатов}
\begin{lmm}\label{lmm:qf-model}
Пусть $\Gamma$~--- полное непротиворечивое множество бескванторных формул. Тогда
существует модель для $\Gamma$.
\end{lmm}
\begin{proof}
Будем строить модель, структурной индукцией по сложности формул. Для начала
разберемся с предметным множеством. В качестве значений для выражений из
констант и функциональных символов зададим строки, содержащие выражения.
Например, $\llbracket c_1\rrbracket = \gn{c_1}$, $\llbracket f_1(c_1,
f_2(c_2))\rrbracket = \gn{f_1(c_1, f_2(c_2))}$ и так далее. Все здесь
происходит аналогично тому, как $\sin(1)$ есть значение, которое мы не можем
вычислить точно и предпочитаем обозначать его своим именем. Таким образом, в
множестве $D$ находятся все возможные выражения, составленные из констант и
функциональных символов.

Теперь рассмотрим формулу~--- некоторый предикат вида $\pi \equiv P(\theta_1,
\dots \theta_n)$, где $\theta_i$~--- это некоторое выражение из функциональных
символов и констант (поскольку все формулы замкнуты, в них не могут участвовать
переменные). Будем считать его истинным, если $\pi \in \Gamma$, иначе, если
$\neg\pi \in \Gamma$, будем считать его ложным.

Все связки получат значения естественным образом.

Теперь покажем, что полученная модель действительно является моделью для
данного множества формул. Возьмем некоторую формулу $\gamma$ из $\Gamma$.
Докажем чуть более сильное свойство: $\llbracket\gamma\rrbracket = \True$ тогда
и только тогда, когда $\gamma \in \Gamma$.

\emph{База.} Очевидно, что если атомарная формула принадлежит $\Gamma$, то она
имеет оценку истина.

\emph{Переход.} Пусть дана некоторая составная формула $\gamma$. Покажем, что
ее оценка истинна тогда и только тогда, когда она входит в $\Gamma$ (при
условии, что это свойство выполнено для составных частей).

Конструкция здесь похожа на конструкцию при доказательстве полноты исчисления
высказываний. Для примера рассмотрим конъюнкцию и отрицание:
\begin{itemize}
\item Пусть $\llbracket\alpha \amp \beta\rrbracket = \True$. Покажем, что
  $\alpha \amp \beta \in \Gamma$.

  В самом деле, пусть это не так и $\neg(\alpha \amp \beta) \in \Gamma$, а,
  значит, $\Gamma, \alpha \amp \beta \proves \psi \amp \neg\psi$. Тогда
  $\Gamma, \alpha, \beta \proves \psi \amp \neg\psi$ (доказуемое утверждение
  $(\psi \amp \varphi \to \pi) \to (\psi \to \varphi \to \pi)$ и теорема о
  дедукции).

  Из таблицы истинности конъюнкции следует, что она истинна только если обе ее
  составных части истинны. То есть $\llbracket\alpha\rrbracket = \True$ и
  $\llbracket\beta\rrbracket = \True$. Значит, $\alpha \in \Gamma$ и $\beta \in
  \Gamma$, что приводит к противоречивости $\Gamma$.
\item Пусть $\llbracket\alpha \amp \beta\rrbracket = \False$. Покажем, что
  $\neg(\alpha \amp \beta) \in \Gamma$.

  Из таблицы истинности следует, что один из параметров обязательно ложен. Пусть,
  например, $\llbracket\alpha\rrbracket = \False$ (случай
  $\llbracket\beta\rrbracket = \False$ рассматривается аналогично). Тогда
  $\alpha \notin \Gamma$, и поэтому $\neg\alpha \in \Gamma$. Рассмотрим
  доказательство:
  \begin{align*}
  &(1)\ \neg\alpha &&\text{Предположение}\\
  &(2)\ \neg\alpha \to \alpha \amp \beta \to \neg\alpha &&\text{Сх. акс. 1}\\
  &(3)\ \alpha \amp \beta \to \neg\alpha &&\text{M.P. 1,2}\\
  &(4)\ \alpha \amp \beta \to \alpha &&\text{Сх. акс. 4}\\
  &(5)\ (\alpha \amp \beta \to \alpha) \to (\alpha \amp \beta \to \neg\alpha) \to \neg(\alpha \amp \beta) &&\text{Сх. акс. 9}\\
  &(6)\ (\alpha \amp \beta \to \neg\alpha) \to \neg(\alpha \amp \beta) &&\text{M.P. 4,5}\\
  &(7)\ \neg(\alpha \amp \beta) &&\text{M.P. 3,6}
  \end{align*}
  Значит, невозможно, чтобы $\alpha \amp \beta \in \Gamma$, поскольку иначе
  получится, что $\Gamma$ противоречиво.
\item Пусть $\llbracket\neg\alpha\rrbracket = \True$. Из оценки следует
  $\llbracket\alpha\rrbracket = \False$, то есть $\neg\alpha \in \Gamma$.
\item Пусть $\llbracket\neg\alpha\rrbracket = \False$. Тогда
  $\llbracket\alpha\rrbracket = \True$. То есть $\alpha \in \Gamma$. Значит,
  невозможно $\neg\alpha \in \Gamma$, иначе $\Gamma$ было бы противоречиво.
\end{itemize}
\end{proof}

\subsection{Теорема Гёделя о полноте исчисления предикатов}
\begin{dfn}
Назовём формулу $\alpha$ формулой с поверхностными кванторами, если существует
такой узел в дереве разбора формулы, не являющийся квантором, ниже которого нет
ни одного квантора, а выше~--- нет ничего, кроме кванторов.
\end{dfn}

Например, формулы $\forall x\exists y\forall z(P(x, y, z) \amp P(z, y, x))$ и
$A \amp B$~--- это формулы с поверхностными кванторами, а формулы $A \amp
\forall x B(x)$, $\exists a P(a) \vee \exists b P(b)$ или $\neg\forall x(P(x)
\to P(y))$ формулами с поверхностными кванторами не являются.

\begin{lmm}
Для любой формулы исчисления предикатов найдётся эквивалентная ей формула с
поверхностными кванторами.
\end{lmm}
\begin{proof}
Доказательство индукцией по структуре формулы. Будем пошагово переносить
кванторы на один уровень выше, при необходимости переименовывая переменные
(если формула имеет вид $\exists x\alpha \amp \exists x\beta$, мы ее в итоге
преобразуем к виду $\exists x_1\exists x_2(\alpha[x := x_1] \amp \beta[x :=
x_2])$) и раскрывая отрицания ($\neg\exists x\alpha$ превратится в $\forall
x\neg\alpha$).
\end{proof}

\begin{thm}[Теорема Гёделя о полноте исчисления предикатов]\label{thm:godel-complete}
Пусть $\Gamma$~--- непротиворечивое множество формул исчисления предикатов.
Тогда существует модель для $\Gamma$.
\end{thm}
\begin{proof}
Без потери общности мы будем предполагать, что $\Gamma$ содержит только формулы
с поверхностными кванторами.

Для доказательства теоремы нам достаточно избавиться от кванторов, не потеряв
непротиворечивости, и сослаться на лемму~\ref{lmm:qf-model}. Для этого мы
определим следующий процесс избавления от одного квантора. Мы построим новый
язык, отличающийся от исходного дополнительными константами. Пусть эти новые
константы имеют имена $d^j_i$ (верхний индекс означает поколение~--- мы сперва
добавим константы $d^1_i$, потом $d^2_i$ и т.п.).

Возьмем непротиворечивое множество формул $\Gamma_g$ и пополним его
дополнительными формулами, получив непротиворечивое множество $\Gamma_{g+1}$,
такое что $\Gamma_g \subseteq \Gamma_{g+1}$.

Возьмем формулу $\gamma \in \Gamma_g$. Возможны такие варианты:
\begin{itemize}
\item $\gamma$ не содержит кванторов. Оставим ее как есть.
\item $\gamma \equiv \forall x\alpha$. Возьмем все константы, использующиеся в
  $\Gamma_g$ (это будут $c_i$ и только те $d^j_i$, которые добавлены раньше,
  т.е. $j \le g$), и все выражения, которые мы можем из них построить с
  участием функциональных символов (их всё равно счётное количество),
  занумеруем их ($\theta_1, \theta_2, \dots$) и добавим формулы $\alpha_1 \equiv
  \alpha[x := \theta_1]$, $\alpha_2 \equiv \alpha[x := \theta_2]$, $\dots$ к
  $\Gamma_{g+1}$.
\item $\gamma \equiv \exists x\alpha$. Возьмем новую константу $d^{g+1}_k$, не
  использовавшуюся ранее, и добавим формулу $\alpha[x := d^{g+1}_k]$ к
  $\Gamma_{g+1}$.
\end{itemize}

Заметим, что все формулы с кванторами пока остаются в $\Gamma_{g+1}$, мы только
добавляем бескванторные формулы. Такая схема позволяет сделать процесс
управляемым (мы всегда добавляем новые переменные, что упрощает доказательства),
но при этом корректным: ведь нам нужно добавлять по формуле $\forall x\alpha$
все возможные выражения $\alpha[x := \theta]$, не только с упоминаемыми в
$\Gamma_g$ константами. При данной схеме мы вернёмся неизбежно к новым (только
что добавленным) константам при следующей итерации и добавим недостающие
формулы.

Покажем, что так построенное множество формул останется непротиворечивым. Пусть
это не так, и существует доказательство противоречия $\Gamma_{g+1} \proves \beta
\amp \neg\beta$. Это доказательство использует конечное количество шагов, и,
следовательно, мы можем явно выписать все его новые по сравнению с $\Gamma$
посылки, перенеся их в правую часть по теореме о дедукции: $\Gamma_g \proves
\gamma_1 \to \dots \gamma_n \to \beta \amp \neg\beta$.

Раз мы оставили справа только новые для $\Gamma_g$ формулы, то каждая из формул
$\gamma_i$ была получена из какой-то исходной формулы из $\Gamma_g$ путем
удаления одного квантора. Мы будем, последовательно перебирая формулы и
перестраивая доказательство, исключать формулы из правой части, пока не окажется,
что предположив противоречивость $\Gamma_{g+1}$, мы получим противоречивость
$\Gamma_g$.

Возможны два варианта:
\begin{itemize}
\item $\gamma_1 \equiv \alpha[x := \theta]$ получено из $\forall x\alpha$. Тогда
  рассмотрим доказательство:
  \begin{align*}
  &(1)\ \forall x\alpha \to \alpha[x := \theta] &&\text{Сх. акс. } \forall\\
  &(2)\ \forall x\alpha &&\forall x\alpha \text{ из } \Gamma_g\\
  &(3)\ \alpha[x := \theta] &&\text{M.P. } 2, 1\\
  &(4\dots k)\ \alpha[x := \theta] \to (\gamma_2 \to \dots \gamma_n \to \beta \amp \neg\beta) &&\text{Исх. формула}\\
  &(k+1)\ \gamma_2 \to \dots \gamma_n \to \beta \amp \neg\beta &&\text{M.P. } 3, k
  \end{align*}
\item $\gamma_1 \equiv \alpha[x := d^{g+1}_k]$ получено из $\exists x\alpha$.
  Выберем какую-нибудь переменную, которая не участвует в выводе противоречия~---
  пусть это $y$. Заменим все вхождения $d^{g+1}_k$ в доказательстве на $y$.
  Поскольку $d^{g+1}_k$~--- константа, то никаких правил для кванторов мы не
  заденем, и доказательство останется верным. Заметим, что поскольку
  $d^{g+1}_k$~--- константа, введенная специально для замены переменной $x$ в
  данной формуле на шаге $g$ и ранее не встречавшаяся~--- то она отсутствует в
  формулах $\gamma_2 \dots \gamma_n$. Также, мы можем правильно выбрать $\beta$,
  чтобы и в нем отсутствовала $d^{g+1}_k$. Значит, мы можем применить правило
  для введения $\exists$:
  \begin{align*}
  &(1\dots k)\ \alpha[x := y] \to (\gamma_2 \to \dots \gamma_n \to \beta \amp \neg\beta) &&\text{Исх. формула}\\
  &(k+1)\ \exists y\,\alpha[x := y] \to (\gamma_2 \to \dots \gamma_n \to \beta \amp \neg\beta) &&\text{Правило для } \exists\\
  &(k+2)\ \exists x\alpha &&\text{Т.к. } \exists x\alpha \text{ из } \Gamma_g\\
  &(k+3\dots l)\ \exists y\,\alpha[x := y] &&\text{Доказуемо}\\
  &(l+1)\ \gamma_2 \to \dots \gamma_n \to \beta \amp \neg\beta &&\text{M.P. } l, k+1
  \end{align*}
\end{itemize}

Взяв $\Gamma_0 \equiv \Gamma$, получим последовательность $\Gamma_0 \subseteq
\Gamma_1 \subseteq \Gamma_2 \subseteq \dots$. Положим $\Gamma^* \equiv \bigcup
\Gamma_i$. $\Gamma^*$ также не может быть противоречиво, поскольку доказательство
использует конечное количество предположений, добавленных, максимум, на шаге
$g$. Значит, множество $\Gamma_g$ тоже противоречиво, что невозможно. Выделив в
$\Gamma^*$ бескванторное подмножество, пополнив его по теореме~\ref{thm:complete-set},
по лемме~\ref{lmm:qf-model} мы получаем модель для него.

Теперь покажем, что это модель и для всего $\Gamma^*$ (а, значит, и для
$\Gamma$).

Рассмотрим некоторую формулу $\gamma \in \Gamma^*$, покажем, что
$\llbracket\gamma\rrbracket = \True$. Проверку мы будем вести индукцией по
структуре «кванторного» префикса формулы.

\emph{База.} Формула не содержит кванторов. В этом случае истинность гарантирует
лемма~\ref{lmm:qf-model}.

\emph{Переход.} Пусть это модель для любой формулы из $\Gamma^*$ с $r$
кванторами. Покажем, что она остаётся моделью и для формул из $\Gamma^*$ с $r+1$
квантором. Пусть формула $\gamma$ впервые добавлена на шаге $p$ к $\Gamma_p$.
Тогда рассмотрим случаи:
\begin{itemize}
\item $\gamma \equiv \forall x\psi$. Нам нужно показать, что формула истинна при
  любом $t \in D$. Возьмём некоторый $t$. По построению модели, есть такое
  $\theta$~--- выражение из констант и функциональных символов, что $\gn{\theta}
  \equiv t$.

  По построению же $\Gamma^*$, начиная с шага $p+1$ мы будем добавлять все
  формулы вида $\psi[x := \kappa]$, где $\kappa$~--- некоторая конструкция из
  констант и функциональных символов. Также, каждая из констант $c_i$ или
  $d^j_i$ из $\theta$ добавлена на некотором шаге $s_k$. То есть, как только и
  константы и формула окажутся в $\Gamma_l$ (понятно, что $l = \max(\max(s_k),
  p)$), так сразу, начиная с $\Gamma_{l+1}$ в нём будет присутствовать и
  $\psi[x := \theta]$. В формуле $\psi$ на один квантор меньше~--- и, по
  предположению индукции, поэтому она истинна.
\item $\gamma \equiv \exists x\psi$. Аналогично, по построению $\Gamma^*$, как
  только $\psi$ добавляется к $\Gamma_g$, так сразу формула $\psi[x :=
  d^{g+1}_k]$ появляется в $\Gamma_{g+1}$. Значит, $\psi$ истинна на значении
  $\gn{d^{g+1}_k}$, то есть $\gamma$ истинна.
\end{itemize}
\end{proof}

\begin{thm}
Если $\models \alpha$, то $\proves \alpha$.
\end{thm}
\begin{proof}
Рассмотрим множество $\Gamma \equiv \{\neg\alpha\}$. Если $\alpha$ недоказуемо,
то $\Gamma$ непротиворечиво, и у $\Gamma$ есть модель, причём $\Gamma \models
\neg\alpha$. Значит, неверно, что $\models \alpha$.
\end{proof}

\section{Теории первого порядка}
Мы занимались до этого момента только логическими рассуждениями самими по себе.
Это интересно, но не очень практически полезно: мы все-таки используем
логические рассуждения для доказательства утверждений о каких-то объектах. Было
бы разумно каким-то образом включить эти объекты в рамки формальной теории.

Одной из основных теорий в данном курсе является формальная арифметика. Но
прежде чем давать формальные определения, поясним ее суть более наглядно.

Рассмотрим некоторое множество $N$. Будем говорить, что оно удовлетворяет
аксиомам Пеано, если выполнено следующее:
\begin{itemize}
\item В нем существует некоторый выделенный элемент $0$.
\item Для каждого элемента множества определена операция $'$, результат ее также
  принадлежит множеству $N$.
\end{itemize}
Кроме того, эти элемент и операция должны удовлетворять следующим требованиям:
\begin{itemize}
\item Не существует такого $x$ из $N$, что $x' = 0$.
\item Если при $x$ и $y$ из $N$ верно $x' = y'$, то $x = y$.
\item Если некоторое утверждение, зависящее от некоторого параметра из $N$,
  верно для $0$, и если из допущения его для $n$ можно вывести его истинность
  для $n'$, то утверждение верно для любого элемента из $N$.
\end{itemize}

Данная аксиоматика позволяет определить натуральные числа (множество
натуральных чисел~--- это множество, удовлетворяющее аксиомам Пеано; заметим,
что тут натуральные числа содержат $0$, так оказывается удобнее) и операции над
ними. Например, сложение можно задать следующими уравнениями:
\begin{align*}
a + 0 &= a\\
a + b' &= (a + b)'
\end{align*}

\begin{thm}
Так определенное сложение коммутативно.
\end{thm}
\begin{proof}
Упражнение.
\end{proof}

Но данная аксиоматика сформулирована неформально, поэтому мы не сможем доказать
никаких содержательных утверждений про нее, пользуясь формальными средствами.
Поэтому мы построим формальную арифметику, теорию, базирующуюся на исчислении
предикатов, в которой также формализована и аксиоматика Пеано.

\begin{dfn}
Теорией первого порядка мы назовем исчисление предикатов с расширенным языком (к
стандартным конструкциям мы добавляем функции и предикаты, возможно, с особым
синтаксисом) и дополнительными аксиомами. Эти дополнения к исчислению
предикатов мы назовем нелогическими или математическими. В противоположность им,
стандартные конструкции и аксиомы исчисления предикатов мы назовем логическими.
\end{dfn}

Возьмем язык исчисления предикатов со следующими изменениями и особенностями:
\begin{itemize}
\item Маленькими латинскими буквами $a, b, \dots$ (возможно, с индексами) будем
  обозначать индивидные переменные.
\item К логическим связкам добавляются такие нелогические конструкции языка:
  $(=)$~--- двуместный предикат, $(+)$ и $(\cdot)$~--- двуместные функции, и
  $(')$~--- одноместная функция. Все двуместные новые конструкции
  левоассоциативны, приоритеты в порядке убывания: $(')$, потом $(\cdot)$,
  потом $(+)$. Все логические связки имеют приоритет ниже. Например,
  $a = b' + b' + c \cdot c \amp b = c$ надо интерпретировать как $(a = (((b') +
  (b')) + (c \cdot c))) \amp (b = c)$.
\item Вводится 0-местная функция $0$.
\end{itemize}

К стандартным аксиомам исчисления предикатов добавим следующие 8 аксиом и одну
схему аксиом.
\begin{align*}
&(A1)\ a = b \to a' = b'\\
&(A2)\ a = b \to a = c \to b = c\\
&(A3)\ a' = b' \to a = b\\
&(A4)\ \neg a' = 0\\
&(A5)\ a + b' = (a + b)'\\
&(A6)\ a + 0 = a\\
&(A7)\ a \cdot 0 = 0\\
&(A8)\ a \cdot b' = a \cdot b + a\\
&(A9)\ (\psi[x := 0]) \amp \forall x((\psi) \to (\psi)[x := x']) \to (\psi)
\end{align*}
В схеме аксиом (A9) $\psi$~--- некоторая формула исчисления предикатов и $x$~---
некоторая переменная, входящая свободно в $\psi$. Данное ограничение объясняется
желанием избежать избыточности, поскольку если $\psi$ не зависит от переменной
$x$, то тогда выражение $\psi \amp \forall x(\psi \to \psi) \to \psi$ является
частным случаем схемы аксиом (4).

\begin{thm}
$\proves a = a$.
\end{thm}
\begin{proof}
Упражнение. Клини, стр.~254.
\end{proof}

\begin{dfn}[Структура]
Структурой теории первого порядка мы назовем упорядоченную тройку $\langle D, F,
P\rangle$, где $F = \langle F_0, F_1, \dots\rangle$~--- списки оценок для
0-местных, 1-местных и т.д. функций, и $P = \langle P_0, P_1, \dots\rangle$~---
списки оценок для 0-местных, 1-местных и т.д. предикатов, $D$~--- предметное
множество.
\end{dfn}

Понятие структуры~--- развитие понятия оценки из исчисления предикатов. Но оно
касается только нелогических составляющих теории; истинностные значения и оценки
для связок по-прежнему определяются исчислением предикатов, лежащим в основе
теории. Для получения оценки формулы нам нужно задать структуру, значения всех
свободных индивидных переменных, и (естественным образом) вычислить результат.

\begin{dfn}
Назовем структуру корректной, если любая доказуемая формула истинна в данной
структуре.
\end{dfn}

\begin{dfn}
Моделью теории мы назовем любую корректную структуру.
\end{dfn}

Еще одним примером теории первого порядка может являться теория групп. К
исчислению предикатов добавим двуместный предикат $(=)$, двуместную функцию
$(\cdot)$, одноместную функцию $(x^{-1})$, константу (т.е. 0-местную функцию)
$1$ и следующие аксиомы:
\begin{align*}
&(E1)\ a = b \to (a = c \to b = c)\\
&(E2)\ a = b \to (a \cdot c = b \cdot c)\\
&(E3)\ a = b \to (c \cdot a = c \cdot b)\\
&(G1)\ a \cdot (b \cdot c) = (a \cdot b) \cdot c\\
&(G2)\ a \cdot 1 = a\\
&(G3)\ a \cdot a^{-1} = 1
\end{align*}

\begin{thm}
Доказуемо, что $a = b \to b = a$ и что $a^{-1} \cdot a = 1$.
\end{thm}
\begin{proof}
Упражнение.
\end{proof}

\section{Рекурсивные функции и отношения}
Рассмотрим примитивы, из которых будем собирать выражения:
\begin{enumerate}
\item $Z : N \to N$, $Z(x) = 0$
\item $N : N \to N$, $N(x) = x'$
\item Проекция. $U^n_i : N^n \to N$, $U^n_i(x_1, \dots x_n) = x_i$
\item Подстановка. Если $f : N^n \to N$ и $g_1, \dots g_n : N^m \to N$, то
  $\Sh{f, g_1, \dots g_n} : N^m \to N$. При этом
  $\Sh{f, g_1, \dots g_n}(x_1, \dots x_m) = f(g_1(x_1, \dots x_m), \dots
  g_n(x_1, \dots x_m))$
\item Примитивная рекурсия. Если $f : N^n \to N$ и $g : N^{n+2} \to N$, то
  $\Rr{f, g} : N^{n+1} \to N$, при этом
  \[
  \Rr{f, g}(x_1, \dots x_n, y) =
  \begin{cases}
  f(x_1, \dots x_n), & y = 0\\
  g(x_1, \dots x_n, y-1, \Rr{f, g}(x_1, \dots x_n, y-1)), & y > 0
  \end{cases}
  \]
\item Минимизация. Если $f : N^{n+1} \to N$, то $\Mn{f} : N^n \to N$, при этом
  $\Mn{f}(x_1, \dots x_n)$~--- такое минимальное число $y$, что $f(x_1, \dots
  x_n, y) = 0$. Если такого $y$ нет, результат данного примитива неопределен.
\end{enumerate}

Если некоторая функция $N^n \to N$ может быть задана с помощью данных
примитивов, то она называется рекурсивной. Если некоторую функцию можно собрать
исключительно из первых 5 примитивов (то есть без использования операции
минимизации), то такая функция называется примитивно-рекурсивной.

\begin{thm}
Следующие функции являются примитивно-рекурсивными: сложение, умножение,
ограниченное вычитание (которое равно 0, если результат вычитания
отрицателен), целочисленное деление, остаток от деления, проверка значения на
простоту.
\end{thm}
\begin{proof}
Упражнение.
\end{proof}

\begin{dfn}
Если $R$~--- $n$-местное отношение, то его характеристической функцией мы
назовем функцию
\[
C_R(x_1, \dots x_n) =
\begin{cases}
1, & \text{если } (x_1, \dots x_n) \in R\\
0, & \text{если } (x_1, \dots x_n) \notin R
\end{cases}
\]
\end{dfn}

\begin{dfn}
Рекурсивным отношением называется отношение, характеристическая функция которого
рекурсивна.
\end{dfn}

\subsection{Рекурсивные, но не примитивно-рекурсивные функции}
Интересен вопрос~--- если столь сложные функции, как вычисление простых чисел,
являются примитивно-рекурсивными, то не достаточно ли их будет для вычисления
любой функции? А если нет, то каков пример функции, не являющейся
примитивно-рекурсивной.

Классическим примером рекурсивной, но не примитивно-рекурсивной функции является
функция Аккермана.

\begin{dfn}
Функцией Аккермана мы назовем так определенную функцию:
\[
A(m, n) =
\begin{cases}
n + 1, & \text{если } m = 0\\
A(m-1, 1), & \text{если } m > 0,\ n = 0\\
A(m-1, A(m, n-1)), & \text{если } m > 0,\ n > 0
\end{cases}
\]
\end{dfn}

Также нам будет удобно другое обозначение: $\alpha_m(n) = A(m, n)$. Данное
обозначение позволяет рассматривать функцию Аккермана как семейство функций от
одной переменной.

\begin{lmm}
$\alpha_{m+1}(n) = \alpha_m^{n+1}(1)$.
\end{lmm}
\begin{proof}
Докажем индукцией по $n$.

\emph{База индукции:} $\alpha_{m+1}(0) = A(m+1, 0) = A(m, 1) = \alpha_m^{0+1}(1)$.

\emph{Переход:} пусть $\alpha_{m+1}(n) = \alpha_m^{n+1}(1)$. Тогда
$\alpha_{m+1}(n+1) = A(m+1, n+1) = A(m, A(m+1, n)) = \alpha_m(\alpha_m^{n+1}(1))
= \alpha_m^{n+2}(1)$.
\end{proof}

\begin{lmm}\label{lmm:ackermann}
Для функции Аккермана справедливы следующие свойства:
\begin{enumerate}
\item Если $m_1 < m_2$, то $A(m_1, n) < A(m_2, n)$
\item Если $n_1 < n_2$, то $A(m, n_1) < A(m, n_2)$
\item $\alpha_1(n) = n + 2$
\item $\alpha_2(n) = 2n + 3$
\item $\alpha_{m+2}(n) > \alpha_m^{n+2}(n)$
\end{enumerate}
\end{lmm}
\begin{proof}
Первые четыре свойства достаточно легко показать по индукции. Покажем
последнее.

Согласно определению функции Аккермана, и предыдущим пунктам леммы, справедлива
цепочка равенств и неравенств:
\begin{multline*}
\alpha_{m+2}(n) = \alpha_{m+1}^{n+1}(1) = \alpha_{m+1}(\alpha_{m+1}^n(1))
= \alpha_m^{\alpha_{m+1}^n(1)+1}(1) = \alpha_m^{\alpha_{m+2}(n-1)+1}(1)\\
\ge \alpha_m^{\alpha_2(n-1)+1}(1) = \alpha_m^{2(n-1)+3+1}(1)
= \alpha_m^{n+n+2}(1) = \alpha_m^{n+2}(\alpha_m^n(1)).
\end{multline*}
Поскольку $\alpha_m^n(1) \ge \alpha_0^n(1) = \mathrm{inc}^n(1) = n+1$, то
$\alpha_m^n(1) > n$, то есть $\alpha_m^{n+2}(\alpha_m^n(1)) > \alpha_m^{n+2}(n)$.
\end{proof}

Введем обозначение: если задан некоторый набор значений $x_1, \dots x_n$, то
для упрощения записи будем писать $\vect{x}$ вместо него. Например, если задана
некоторая функция $f : N^n \to N$, то будем писать $f(\vect{x})$ вместо $f(x_1,
\dots, x_n)$. Запись же $g(\vect{x}, y)$ означает, что функция $g$ применяется к
$n+1$ аргументу: $g(x_1, \dots, x_n, y)$.

\begin{thm}\label{thm:ack-grows}
Функция Аккермана растет быстрее любой примитивно-рекурсивной функции. Точнее,
какова бы ни была примитивно-рекурсивная функция $p : N^n \to N$, мы можем
подобрать такую константу $K$, что $p(\vect{x}) \le \alpha_K(\max(\vect{x}))$
при любом $\vect{x}$.
\end{thm}
\begin{proof}
Докажем индукцией по структуре формулы, показывающей примитивную рекурсивность
$p$. Базой будут являться примитивы $Z$, $N$ и $U$, для которых утверждение
очевидно. Покажем переход для примитивов $S$ и $R$.

Рассмотрим примитив $S$. Пусть $p = \Sh{f, g_1, \dots g_m}$, где $f, g_1, \dots
g_m$~--- это некоторые примитивно-рекурсивные функции, для которых по
предположению индукции уже найдены требуемые константы $K_0, K_1 \dots K_m$.
Тогда возьмем $K = \max(K_0 \dots K_m) + 2$, и покажем, что это значение нас
устроит. В самом деле:
\[
p(\vect{x}) = f(g_1(\vect{x}), \dots g_m(\vect{x}))
\le \alpha_{K_0}(\max(g_1(\vect{x}), \dots g_m(\vect{x})))
\le \alpha_{K-2}(\max(g_1(\vect{x}), \dots g_m(\vect{x})))
\]
В свою очередь, каждый $g_i$ мы можем оценить через соответствующий
$\alpha_{K_i}$:
\[
g_i(\vect{x}) \le \alpha_{K_i}(\max(\vect{x})) \le \alpha_{K-2}(\max(\vect{x}))
\]
из чего получается, что
\[
\alpha_{K-2}(\max(g_1(\vect{x}), \dots g_m(\vect{x})))
\le \alpha_{K-2}(\alpha_{K-2}(\max(\vect{x})))
\le \alpha_{K-2}^{\max(\vect{x})+2}(\max(\vect{x}))
\]
Применив свойство из предыдущей леммы, мы можем заключить требуемое:
\[
\alpha_{K-2}^{\max(\vect{x})+2}(\max(\vect{x})) < \alpha_K(\max(\vect{x}))
\]

Теперь рассмотрим примитив $R$: $p = \Rr{f, g}$, причём $f : N^{n-1} \to N$, и
$g : N^{n+1} \to N$, и для них уже найдены константы $K_1$ и $K_2$,
соответственно. Тогда возьмем $K = \max(K_1, K_2) + 2$.

Для доказательства утверждения покажем чуть более сильное утверждение:
\[
p(\vect{x}, y) \le \alpha_{K-2}^{y+1}(\max(\vect{x}, y))
\]
Этого достаточно, поскольку
\[
\alpha_{K-2}^{y+1}(\max(\vect{x}, y)) \le \alpha_{K-2}^{\max(\vect{x},y)+1}(\max(\vect{x}, y)) \le
\alpha_{K-2}^{\max(\vect{x},y)+2}(\max(\vect{x}, y)) < \alpha_K(\max(\vect{x}, y))
\]
Покажем требуемое индукцией по $y$. База: $y = 0$. Тогда:
\[
p(\vect{x}, 0) = f(\vect{x}) \le \alpha_{K_1}(\max(\vect{x})) = \alpha_{K_1}(\max(\vect{x}, 0))
\le \alpha_{K-2}(\max(\vect{x}, 0))
\]
Переход: пусть $p(\vect{x}, y) \le \alpha_{K-2}^{y+1}(\max(\vect{x}, y))$. Тогда
\begin{multline*}
p(\vect{x}, y+1) = g(\vect{x}, y, p(\vect{x}, y))
\le \alpha_{K_2}(\max(\vect{x}, y, p(\vect{x}, y)))\\
\le \alpha_{K_2}(\max(\vect{x}, y, \alpha_{K-2}^{y+1}(\max(\vect{x}, y)))) =\\
\text{(тут используется монотонность } \alpha) \\
= \alpha_{K_2}(\alpha_{K-2}^{y+1}(\max(\vect{x}, y)))
\le \alpha_{K-2}^{y+2}(\max(\vect{x}, y+1))
\end{multline*}
\end{proof}

\begin{thm}
Функция Аккермана не является примитивно-рекурсивной.
\end{thm}
\begin{proof}
Пусть это не так, и функция $A(x, y)$~--- примитивно-рекурсивна. Тогда
рассмотрим $g(x) = A(x, x)$. Эта функция, очевидно, тоже является
примитивно-рекурсивной.

Однако, по предыдущей теореме, найдется такой $N$, что $g(x) \le A(N, x)$ при
любом $x$. В том числе это верно и при $x = N+1$. То есть, по предыдущей
теореме, $g(N+1) = A(N+1, N+1) \le A(N, N+1)$, однако нам также известно, что
$A(N+1, N+1) > A(N, N+1)$.
\end{proof}

\section{Арифметические отношения и функции. Их выразимость и представимость в формальной арифметике}
\subsection{Выразимость и представимость}
Введем обозначение. Если в тексте вводится некоторая формула $\alpha(x_1, \dots
x_n)$, то по умолчанию считается, что эта формула имеет минимум $n$ свободных
переменных, с именами $x_1, \dots x_n$.

Внутри же выражения запись $\alpha(y_1, \dots y_n)$ мы будем трактовать, как
$\alpha[x_1 := y_1, \dots x_n := y_n]$, при этом мы подразумеваем, что $y_1,
\dots y_n$ свободны для подстановки вместо $x_1, \dots x_n$ в $\alpha$.

Также, запись $B(x_1, \dots x_n) \equiv \alpha(x_1, \dots x_n)$ будет означать,
что мы определяем новую формулу с именем $B$ и $n$ свободными переменными $x_1,
\dots x_n$. Данная формула должна восприниматься только как сокращение записи,
макроподстановка.

\begin{dfn}
Арифметическая функция~--- функция $f : N^n \to N$. Арифметическое отношение~---
$n$-арное отношение, заданное на $N$.
\end{dfn}

\begin{dfn}
Арифметическое отношение $R$ называется выразимым (в формальной арифметике),
если существует такая формула $\alpha(x_1, \dots x_n)$ с $n$ свободными
переменными, что для любых натуральных чисел $k_1 \dots k_n$
\begin{enumerate}
\item если $(k_1, \dots k_n) \in R$, то доказуемо $\alpha(k_1, \dots k_n)$
\item если $(k_1, \dots k_n) \notin R$, то доказуемо $\neg\alpha(k_1, \dots
  k_n)$.
\end{enumerate}
\end{dfn}

Например, отношение $(<)$ является выразимым в арифметике: рассмотрим формулу
$\alpha(a_1, a_2) = \exists b(\neg b = 0 \amp a_1 + b = a_2)$. В самом деле, если
взять некоторые числа $k_1$ и $k_2$, такие, что $k_1 < k_2$, то найдется такое
положительное число $b$, что $k_1 + b = k_2$. Можно показать, что если
подставить $k_1$ и $k_2$ в $\alpha$, то формула будет доказуема.

Наметим доказательство: тут должно быть два доказательства по индукции, сперва
по $k_2$, потом по $k_1$. Рассмотрим доказательство по индукции: пусть $k_1 =
0$, индукция по 2-му параметру: разберем доказательство базы при $k_2 = 1$.
Тогда надо показать $\exists b(\neg b = 0 \amp 0 + b = 1)$:
\begin{align*}
&(1)\ \neg 1 = 0 \amp 0 + 1 = 1 &&\text{Несложно показать}\\
&(2)\ (\neg 1 = 0 \amp 0 + 1 = 1) \to \exists b(\neg b = 0 \amp 0 + b = 1) &&\text{Сх. акс. для } \exists\\
&(3)\ \exists b(\neg b = 0 \amp 0 + b = 1) &&\text{M.P. 1 и 2.}
\end{align*}

\begin{dfn}
Введем следующее сокращение записи: пусть $\exists!y\,\varphi(y)$ означает
\[ \exists y\,\varphi(y) \amp \forall a\forall b(\varphi(a) \amp \varphi(b) \to a = b) \]
Здесь $a$ и $b$~--- некоторые переменные, не входящие в формулу $\varphi$
свободно.
\end{dfn}

\begin{dfn}
Арифметическая функция $f$ от $n$ аргументов называется представимой в
формальной арифметике, если существует такая формула $\alpha(x_1, \dots
x_{n+1})$ с $n+1$ свободными пременными, что для любых натуральных чисел $k_1
\dots k_{n+1}$
\begin{enumerate}
\item $f(k_1, \dots k_n) = k_{n+1}$ тогда и только тогда, когда доказуемо
  $\alpha(k_1, \dots k_{n+1})$.
\item Доказуемо $\exists!b\,\alpha(k_1, \dots k_n, b)$.
\end{enumerate}
\end{dfn}

\subsection{Представимость рекурсивных функций}
\begin{thm}
Функции $Z$, $N$, $U^n_i$ являются представимыми.
\end{thm}
\begin{proof}
Наметим доказательство. Для этого приведем формулы, доказательство корректности
этих формул оставим в виде упражнения.
\begin{itemize}
\item Примитив $Z$ представит формула $Z(a, b) := (a = a \amp b = 0)$.
\item Примитив $N$ представит формула $N(a, b) := (a' = b)$.
\item Примитив $U^n_i$ представит формула $U^n_i(a_1, \dots a_n, b) = (a_1 = a_1)
  \amp \dots \amp (a_n = a_n) \amp (b = a_i)$.
\end{itemize}
\end{proof}

\begin{thm}
Если функции $f$ и $g_1, \dots g_m$ представимы, то функция $\Sh{f, g_1, \dots
g_m}$ также представима.
\end{thm}
\begin{proof}
Поскольку функции $f$ и $g_i$ представимы, то есть формулы $F$ и $G_1, \dots
G_m$, их представляющие. Тогда следующая формула представит $\Sh{f, g_1, \dots
g_m}$:
\[
S(a_1, \dots a_n, b) := \exists b_1 \dots \exists b_m(G_1(a_1, \dots a_n, b_1)
\amp \dots \amp G_m(a_1, \dots a_n, b_m) \amp F(b_1, \dots b_m, b))
\]
\end{proof}

\begin{dfn}
$\beta$-функция Гёделя~--- это функция $\beta(b, c, i) = b \% (1 + c \cdot (i +
1))$. Здесь операция $(\%)$ означает взятие остатка от целочисленного деления.
\end{dfn}

\begin{lmm}
Функция примитивно-рекурсивна, и при этом представима в арифметике формулой
\[ B(b, c, i, d) := \exists q((b = q \cdot (1 + c \cdot (i + 1)) + d) \amp (d < 1 + c \cdot (i + 1))) \]
\end{lmm}
\begin{proof}
Упражнение.
\end{proof}

\begin{thm}[Китайская теорема об остатках]
Если $u_1, \dots u_n$~--- попарно взаимно простые целые числа, и $k_1, \dots
k_n$~--- целые числа, такие, что $0 \le k_i < u_i$ при любом $i$, то найдется
такое целое число $b$, что $k_i = b \% u_i$ при любом $i$.
\end{thm}
Доказательство этой теоремы по индукции несложно, но не входит в наш курс.

\begin{lmm}
Для любой конечной последовательности чисел $k_0 \dots k_n$ можно подобрать
такие константы $b$ и $c$, что $\beta(b, c, i) = k_i$ для $0 \le i \le n$.
\end{lmm}
\begin{proof}
Возьмем число $c = \max(k_1, \dots k_n, n)!$. Рассмотрим числа $u_i = 1 + c
\cdot (i + 1)$.
\begin{itemize}
\item Никакие числа $u_i$ и $u_j$ ($0 \le j < i \le n$) не имеют общих делителей
  кроме 1. Пусть это не так, и есть некоторый общий делитель $p$ (очевидно, мы
  можем предположить его простоту~--- разложив на множители, если он составной).
  Тогда $p$ будет делить $u_i - u_j = c \cdot (i - j)$, при этом $p$ не может
  делить $c$~--- иначе окажется, что $u_i = (1 + c \cdot (i + 1))$ делится на
  $p$ и $c \cdot (i + 1)$ делится на $p$. Значит, $p$ делит $i - j$, то есть все
  равно делит $c$, так как $c$~--- факториал некоторого числа, не меньшего $n$,
  и при этом $i - j \le n$.
\item Каждое из чисел $k_i$ меньше, чем $u_i$: в самом деле, $k_i \le c < 1 + c
  \cdot (i + 1) = u_i$.
\end{itemize}
Из полученных утверждений видно, что выполнены условия китайской теоремы об
остатках. Следовательно, найдется такое целое число $b$, для которого выполнено
$k_i = b \% u_i$, то есть $k_i = \beta(b, c, i)$.
\end{proof}

\begin{thm}
Всякая рекурсивная функция представима в арифметике.
\end{thm}
\begin{proof}
Представимость первых четырех примитивов уже показана. Покажем представимость
примитивной рекурсии и операции минимизации.

\emph{Примитивная рекурсия.} Пусть есть некоторый $\Rr{f, g}$. Соответственно,
$f$ и $g$ уже представлены как некоторые формулы $F$ и $G$. Из определения
$\Rr{f, g}$ мы знаем, что для значения $\Rr{f, g}(x_1, \dots x_{n+1})$ должна
существовать последовательность $a_0 \dots a_{x_{n+1}}$ результатов применения
функций $f$ и $g$~--- значений на одно больше, чем итераций в цикле примитивной
рекурсии, а это количество определяется последним параметром функции $\Rr{f,
g}$. При этом:
\begin{align*}
a_0 &= f(x_1, \dots x_n)\\
a_1 &= g(x_1, \dots x_n, 0, a_0)\\
&\ \,\vdots\\
a_{x_{n+1}} &= g(x_1, \dots x_n, x_{n+1} - 1, a_{x_{n+1}-1})
\end{align*}
Значит, по лемме, должны существовать такие числа $b$ и $c$, что $\beta(b, c, i)
= a_i$ для $0 \le i \le x_{n+1}$.

Приведенные рассуждения позволяют построить следующую формулу, представляющую
$\Rr{f, g}(x_1, \dots x_{n+1})$:
\begin{multline*}
R(x_1, \dots x_{n+1}, a) := \exists b\exists c(\exists k(B(b, c, 0, k) \amp F(x_1, \dots x_n, k))\\
\amp B(b, c, x_{n+1}, a)\\
\amp \forall k(k < x_{n+1} \to \exists d\exists e(B(b, c, k, d) \amp B(b, c, k', e) \amp G(x_1, \dots x_n, k, d, e))))
\end{multline*}

\emph{Минимизация.} Рассмотрим конструкцию $\Mn{f}$. $f$ уже представлено как
некоторая формула $F$. Тогда формула
\[ M(x_1, \dots x_n, y) := F(x_1, \dots x_n, y, 0) \amp \forall z(z < y \to \neg F(x_1, \dots x_n, z, 0)) \]
представит $\Mn{f}$.
\end{proof}

\begin{thm}
Всякое рекурсивное арифметическое отношение выразимо в формальной арифметике.
\end{thm}
\begin{proof}
Рассмотрим характеристическую функцию $C_R$ некоторого рекурсивного отношения
$R$. Данная функция (как любая другая рекурсивная функция) представима в
формальной арифметике~--- пусть ее представляет формула $\rho(x_1, \dots x_n,
y)$. Тогда формула $\rho(x_1, \dots x_n, 1)$ выразит соответствующее отношение.

В самом деле, если $(k_1, \dots k_n) \in R$ для некоторых $k_i$, то $\proves
\rho(k_1, \dots k_n, 1)$.

Обратно, пусть $(k_1, \dots k_n) \notin R$. Из представимости $C_R$ известно
\[ \proves \exists!y\,\rho(k_1, \dots k_n, y) \]
то есть, в частности,
\[ \proves \forall y_1\forall y_2(\rho(k_1, \dots k_n, y_1) \amp \rho(k_1, \dots k_n, y_2) \to y_1 = y_2) \]
что можно преобразовать в
\[ \proves \rho(k_1, \dots k_n, 0) \to (\rho(k_1, \dots k_n, 1) \to 1 = 0) \]
Поскольку $\proves \rho(k_1, \dots k_n, 0)$ доказуемо вследствие представимости
$C_R$, то мы можем далее вывести
\[ \proves \rho(k_1, \dots k_n, 1) \to 1 = 0 \]
и, заметив, что $\proves \neg 1 = 0$, в итоге покажем требуемое для выразимости
\[ \proves \neg\rho(k_1, \dots k_n, 1) \]
\end{proof}

\section{Гёделева нумерация. Арифметизация доказательств}
Ранее мы показали, что любое рекурсивное арифметическое отношение выразимо в
формальной арифметике. Теперь мы покажем, что любое выразимое в формальной
арифметике отношение является рекурсивным.

\begin{dfn}
Ограниченные кванторы $\exists x{<}y\,\varphi(x)$ и $\forall x{<}y\,\varphi(x)$~---
сокращения записи для выражений вида $\exists x(x < y \amp \varphi(x))$ и
$\forall x(x \ge y \vee \varphi(x))$.
\end{dfn}

\begin{thm}
Пусть $P_1$ и $P_2$~--- рекурсивные отношения. Тогда следующие комбинации
отношений также являются рекурсивными отношениями:
\begin{enumerate}
\item $F(x_1, \dots x_n, z) := \forall y{<}z\,P_1(x_1, \dots x_n, y)$
\item $E(x_1, \dots x_n, z) := \exists y{<}z\,P_1(x_1, \dots x_n, y)$
\item $P_1(x_1, \dots x_n) \to P_2(x_1, \dots x_n)$
\item $P_1(x_1, \dots x_n) \vee P_2(x_1, \dots x_n)$
\item $P_1(x_1, \dots x_n) \amp P_2(x_1, \dots x_n)$
\item $\neg P_1(x_1, \dots x_n)$
\end{enumerate}
\end{thm}
\begin{proof}
Упражнение.
\end{proof}

Теперь мы перенесем понятие вывода формулы на язык рекурсивных отношений, и,
следовательно, внутрь языка формальной арифметики.

\begin{dfn}[Гёделева нумерация]
Дадим следующие номера символам языка формальной арифметики:
\[
\begin{array}{rl}
3 & (\\
5 & )\\
7 & ,\\
9 & \neg\\
11 & \to\\
13 & \vee\\
15 & \amp\\
17 & \forall\\
19 & \exists\\
21 + 6 \cdot k & x_k \text{ переменные}\\
23 + 6 \cdot 2^k \cdot 3^n & f^n_k \text{ $n$-местные функциональные символы: } (') , (+) \text{ и т.п.}\\
25 + 6 \cdot 2^k \cdot 3^n & P^n_k \text{ $n$-местные предикаты, в т.ч. } (=)
\end{array}
\]
\end{dfn}

Уточним язык~--- обяжем всегда писать скобки вокруг двуместной операции, и
только в этом случае. Также укажем номера для предопределенных функций и
предикатов: $f^1_0 \equiv (')$, $f^2_0 \equiv (+)$, $f^2_1 \equiv (\cdot)$,
$P^2_0 \equiv (=)$.

Научимся записывать выражения в виде чисел. Пусть $p_1, \dots p_k, \dots$~---
список простых чисел, при этом $p_1 = 2$, $p_2 = 3$, $\dots$.

Тогда текст из $n$ символов с гёделевыми номерами $c_1, \dots c_n$ запишем как
число $t = p_1^{c_1} \cdot p_2^{c_2} \cdot \dots \cdot p_n^{c_n}$. Ясно, что
такое представление однозначно позволяет установить длину строки (гёделева
нумерация не содержит 0, поэтому можно определить длину строки как максимальный
номер простого числа, на которое делится $t$; будем записывать эту функцию как
$\mathrm{Len}(s)$), и каждый символ строки в отдельности (будем записывать
функцию как $(s)_n$). Также ясно, что функции $\mathrm{Len}$ и $(x)_n$~---
рекурсивны.

Чтобы удобнее работать со строками, введем следующие обозначения, за которыми
скрываются рекурсивные функции:
\begin{itemize}
\item запись $\gn{c_1 c_2 c_3 \dots}$ (где $c_i$~--- какие-то символы языка
  формальной арифметики) задает рекурсивную функцию $f : N \to N$, при этом
  $f(x) = p_1^{c_1} \cdot \dots \cdot p_n^{c_n}$.
\item Операцию конкатенации строк определим так. Пусть $S$ и $T$~--- рекурсивные
  функции, результатом вычисления которых являются числа $s = p_1^{s_1} \cdot
  \dots \cdot p_n^{s_n}$ и $t = p_1^{t_1} \cdot \dots \cdot p_m^{t_m}$
  соответственно. Тогда $S \mathbin{@} T$~--- это рекурсивная функция,
  вычисляющая функции $S$ и $T$, результатом работы которой будет $p_1^{s_1}
  \cdot \dots \cdot p_n^{s_n} \cdot p_{n+1}^{t_1} \cdot \dots \cdot
  p_{n+m}^{t_m}$.
\end{itemize}

Чтобы представить доказательства, мы будем объединять строки вместе так же, как
объединяем символы в строки: $2^{2^3} \cdot 3^{2^5}$~--- это последовательность
из двух строк, первая~--- это «(», а вторая~--- «)».

Теперь мы можем понять, как написать программу, проверяющую корректность
доказательства некоторого утверждения в формальной арифметике. Наметим общую
идею. Программа будет состоять из набора рекурсивных отношений и функций,
каждое из которых выражает некоторое отношение, содержательное для проверки
доказательства. Ниже мы покажем идею данной конструкции, приведя несколько из
них.
\begin{itemize}
\item Проверка того, что $a$~--- строка, состоящая только из переменной.
  $\mathrm{Var}(a) \equiv \exists z{<}a\ (a = 2^{21+6\cdot z})$
\item Проверка того, что выражение с номером $a$ получено из выражений $b$ и $c$
  путем применения правила Modus Ponens. $\mathrm{MP}(b, c, a) \equiv c =
  \gn{(} \mathbin{@} b \mathbin{@} \gn{\to} \mathbin{@} a \mathbin{@} \gn{)}$
\item Проверка того, что $b$ получается из $a$ подстановкой $y$ вместо $x$:
  $\mathrm{Subst}(a, b, x, y)$~--- без реализации
\item Функция, подставляющая $y$ вместо $x$ в формуле $a$:
  $\mathrm{Sub}(a, x, y) \equiv \Mn{\Sh{\mathrm{Subst}, U^4_1, U^4_4, U^4_2,
  U^4_3}}(a, x, y)$
\item Проверка того, что переменная номер $x$ входит свободно в формулу $f$.
  $\mathrm{Free}(f, x) \equiv \neg\mathrm{Subst}(a, a, x, 21 + 6x)$
\item Функция, выдающая гёделев номер выражения, соответствующего целому числу:
  $\mathrm{Num} \equiv \Rr{\gn{0}, \Sh{@, U^3_3, \Sh{\gn{'}, U^3_1}}}$ \ (с
  точностью до вспомогательных проекций $U^1_1$)
\end{itemize}

Путем некоторых усилий мы можем выписать формулу, представляющую двуместное
отношение $\mathrm{Proof}(f, p)$, истинное тогда и только тогда, когда $p$~---
гёделев номер доказательства формулы с гёделевым номером $f$.

\begin{thm}
Любая представимая в формальной арифметике функция является рекурсивной.
\end{thm}
\begin{proof}
Возьмем некоторую представимую функцию $f : N^n \to N$. Значит, для нее
существует формула формальной арифметики, представляющая ее. Пусть $\varphi$~---
эта формула (со свободными переменными $x_1, \dots x_n, y$); при этом в случае
$f(u_1, \dots u_n) = v$ должно быть доказуемо $\varphi(u_1, \dots u_n, v)$. По
формуле можно построить рекурсивную функцию, $C_\varphi(u_1, \dots u_n, v, p)$,
выражающую тот факт, что $p$~--- гёделев номер вывода формулы $\varphi(u_1,
\dots u_n, v)$. Тогда возьмем
\[
f(x_1, \dots x_n) \equiv \bigl(\Mn{\Sh{C_\varphi, U^{n+1}_1, \dots U^{n+1}_n,
(U^{n+1}_{n+1})_1, (U^{n+1}_{n+1})_2}}(x_1, \dots x_n)\bigr)_1.
\]
\end{proof}

\section{1-я и 2-я теоремы Гёделя о неполноте арифметики}
\begin{dfn}
Мы будем называть теорию непротиворечивой, если не найдется такой формулы
$\alpha$, что доказуемо как $\alpha$, так и $\neg\alpha$.
\end{dfn}

\begin{lmm}
Если теория противоречива, то в ней доказуема любая формула.
\end{lmm}
\begin{proof}
Если теория противоречива, то в ней есть утверждение $\alpha$, что доказуемо
$\alpha$ и $\neg\alpha$. Воспользуемся общезначимой (и потому доказуемой)
формулой исчисления высказываний $\alpha \to \neg\alpha \to \beta$.
\end{proof}

\begin{dfn}
Мы будем называть теорию $\omega$-непротиворечивой, если, какова бы ни была
формула $P(x)$ со свободной переменной $x$, такая, что для любого натурального
числа $p$ доказуемо $P(p)$, то формула $\exists p\neg P(p)$ недоказуема.
\end{dfn}

\begin{lmm}
$\omega$-непротиворечивость влечёт непротиворечивость.
\end{lmm}
\begin{proof}
Рассмотрим выводимую формулу $x = x \to x = x$. При подстановке любого
натурального числа вместо $x$ формула будет по-прежнему выводима: $k = k \to
k = k$. Значит, по $\omega$-непротиворечивости формула $\exists p\neg(x = x \to
x = x)$ невыводима. Значит, теория непротиворечива (поскольку в противоречивой
теории выводится любая формула).
\end{proof}

Определим рекурсивное отношение $W_1(x, p)$, истинное тогда и только тогда,
когда $x$ есть гёделев номер некоторой формулы $\varphi$ с единственной
свободной переменной $x$, а $p$ есть гёделев номер доказательства
$\varphi(\gn{\varphi})$~--- доказательства самоприменения $\varphi$. Это
соотношение является выразимым в формальной арифметике. В самом деле, его можно
задать с помощью рекурсивных отношений, введённых ранее:
\[ W_1(x, p) := \mathrm{Free}(x, \gn{x}) \amp \mathrm{Proof}(\mathrm{Sub}(x, \gn{x}, \mathrm{Num}(x)), p) \]
В силу выразимости рекурсивных отношений в формальной арифметике, для неё
найдётся формула $\omega_1(x, p)$, причём, если $\langle x, p\rangle \in W_1$, то
доказуемо $\omega_1(x, p)$, а если $\langle x, p\rangle \notin W_1$, то уже
доказуемо $\neg\omega_1(x, p)$.

Рассмотрим формулу $\sigma \equiv \forall p\neg\omega_1(x, p)$~--- это
некоторая формула с единственной свободной переменной $x$~--- и посмотрим, что
произойдёт с её самоприменением: $\sigma(\gn{\sigma})$. Внимательное
наблюдение за происходящим даст следующую теорему.

\begin{thm}[Первая теорема Гёделя о неполноте арифметики]
\begin{enumerate}
\item Если формальная арифметика непротиворечива, то недоказуемо
  $\sigma(\gn{\sigma})$.
\item Если формальная арифметика $\omega$-непротиворечива, то недоказуемо
  $\neg\sigma(\gn{\sigma})$.
\end{enumerate}
\end{thm}
\begin{proof}
1. Пусть $\proves \sigma(\gn{\sigma})$. Тогда найдется гёделев номер ее
доказательства $p$, значит, $W_1(\gn{\sigma}, p)$, то есть $\proves
\omega_1(\gn{\sigma}, p)$. С другой стороны, пользуясь схемой аксиом для
квантора всеобщности и правилом Modus Ponens, из предположения теоремы $\proves
\sigma(\gn{\sigma})$ (то есть $\proves \forall p\neg\omega_1(\gn{\sigma}, p)$)
можно показать $\proves \neg\omega_1(\gn{\sigma}, p)$. Противоречие.

2. Пусть $\proves \neg\sigma(\gn{\sigma})$, то есть $\proves \neg\forall
p\neg\omega_1(\gn{\sigma}, p)$, то есть $\proves \exists p\,\omega_1(\gn{\sigma},
p)$. Значит, неизбежно найдется такой номер $q$, что $\proves
\omega_1(\gn{\sigma}, q)$, поскольку если бы для каждого $q$ было бы доказуемо
$\proves \neg\omega_1(\gn{\sigma}, q)$, то по $\omega$-непротиворечивости было
бы недоказуемо $\exists p\neg\neg\omega_1(\gn{\sigma}, p)$.

Рассмотрев же определение $W_1$, можно заметить, что найденный $q$ также есть
номер доказательства $\sigma(\gn{\sigma})$, что вступает в противоречие с
предположением $\proves \neg\sigma(\gn{\sigma})$.
\end{proof}

Формула $\sigma(\gn{\sigma})$, говоря простым языком, утверждает собственную
недоказуемость. Мы показали, что эта формула (при условии
$\omega$-непротиворечивости формальной арифметики) действительно недоказуема~---
что означает её общезначимость. Таким образом, мы показали, что если формальная
арифметика $\omega$-непротиворечива, то она неполна.

В данном рассуждении используется сложное понятие $\omega$-непротиворечивости,
что смущает. Теорема Гёделя в форме Россера снимает эту сложность.

Рассмотрим отношение $W_2(x, p)$~--- $x$ и $p$ состоят в отношении $W_2$ тогда и
только тогда, когда $p$~--- гёделев номер доказательства отрицания
самоприменения $x$ (если $\varphi$~--- формула от одной переменной $x$, то
$p$~--- номер доказательства $\neg\varphi(\gn{\varphi})$). Мы также можем
выразить его в формальной арифметике аналогично $\omega_1$ (обозначим выражающую
формулу за $\omega_2$).

Тогда рассмотрим формулу $\rho(x) \equiv \forall p(\omega_1(x, p) \to \exists
q(q < p \amp \omega_2(x, q)))$. Неформальным языком она утверждает, что для
любого доказательства самоприменения некоторой формулы с номером $a$ найдется
доказательство (да еще и с меньшим гёделевым номером) отрицания этой формулы. Ну
и по традиции применим ее к своему номеру $r$. Внимательное рассмотрение этой
ситуации приводит к следующей теореме.

\begin{thm}[Теорема Гёделя в форме Россера]
Если формальная арифметика непротиворечива, то не доказуема как формула
$\rho(\gn{\rho})$, так и её отрицание.
\end{thm}

Докажем эту теорему, рассмотрев вспомогательную лемму:

\begin{lmm}
Каково бы ни было число $n$, доказуемы следующие утверждения:
\begin{itemize}
\item $\proves a \le n \to (a = 0 \vee a = 1 \vee \dots \vee a = n)$
\item $\proves (a = 0 \vee a = 1 \vee \dots \vee a = n) \to a \le n$
\end{itemize}
\end{lmm}
\begin{proof}
Импликации доказываются индукцией по $n$. Мы не будем предлагать подробного
доказательства (в силу его размера и технического характера), наметим только
несколько шагов. Доказательство значительно объемнее, но делается достаточно
похоже.

Докажем первую импликацию. Рассмотрим индукцию по $b$ на мета-языке.

\emph{База.} $n = 0$. Тогда $a \le n$ после подстановки $b$ и раскрытия
определения отношения «меньше» превращается в $\exists b(a + b = 0)$. Нам нужно
показать, что $\proves \exists b(a + b = 0) \to a = 0$.

Утверждение может быть получено применением правила введения $\exists$ из более
простого: $\proves a + b = 0 \to a = 0$.

Докажем данное утверждение индукцией по $b$ в предметном языке~--- применив
схему аксиом индукции.

Рассмотрим сокращение записи: $A(b) \equiv a + b = 0 \to a = 0$, тогда
следующее выражение~--- аксиома (по схеме A9): $A(0) \amp \forall b(A(b) \to
A(b')) \to A(b)$. Если показать $A(0)$ и $\forall b(A(b) \to A(b'))$, то из
этого по правилу M.P. будет нетрудно получить необходимое $A(b)$.

Покажем $A(0)$.
\begin{align*}
&(1\dots l)\ a + 0 = 0 \to a + 0 = a \to a = 0 &&\text{Акс. А2 + переим.}\\
&(l+1)\ a + 0 = a \to a + 0 = 0 \to a + 0 = a &&\text{Сх. акс. 1}\\
&(l+2)\ a + 0 = a &&\text{Акс. А6}\\
&(l+3)\ a + 0 = 0 \to a + 0 = a &&\text{M.P.}\\
&(l+4\dots m)\ a + 0 = 0 \to a = 0 &&\text{Сх. акс. 2 + M.P.}
\end{align*}
Теперь покажем $A(b) \to A(b')$.
\begin{align*}
&(1)\ a + b' = (a + b)' \to a + b' = 0 \to (a + b)' = 0 &&\text{Сх. акс. 1}\\
&(2)\ a + b' = 0 \to (a + b)' = 0 &&\text{Акс. А5}\\
&(3\dots k)\ a + b' = 0 \to \neg(a + b)' = 0 &&\text{Акс. А4 + замена пер. + ослабление}\\
&(k+1)\ (a + b' = 0 \to (a + b)' = 0) \to (a + b' = 0 \to \neg(a + b)' = 0) \to (\neg a + b' = 0) &&\\
&(l)\ \neg a + b' = 0 &&\text{M.P. 2 раза}\\
&(l+1\dots m)\ (a + b' = 0) \to (\neg a + b' = 0) \to (a = 0) &&\text{Инт. сх. акс. 10}\\
&(m+1\dots p)\ (a + b' = 0) \to (a = 0) &&\text{Сх. акс. 2 + M.P. 2 раза}\\
&(p+1\dots q)\ (a + b = 0 \to a = 0) \to (a + b' = 0 \to a = 0) &&\text{Ослабление}
\end{align*}
\end{proof}

\begin{proof}[Доказательство теоремы Гёделя в форме Россера]
Пусть $\proves \rho(\gn{\rho})$, т.е. $\proves \forall p(\omega_1(\gn{\rho}, p)
\to \exists q(q < p \amp \omega_2(\gn{\rho}, q)))$. У этого доказательства есть
некоторый номер $t$, причем $\omega_1(\gn{\rho}, t)$. Несложно показать $\proves
(\omega_1(\gn{\rho}, t) \to \exists q(q < t \amp \omega_2(\gn{\rho}, q)))$. По
построению формулы $\proves \omega_1(\gn{\rho}, t)$, откуда, соединив
доказательства воедино, и применив правило Modus Ponens, получим $\proves
\exists q(q < t \amp \omega_2(\gn{\rho}, q))$.

Так как теория непротиворечива, то не существует вывода формулы $\forall
p(\omega_1(\gn{\rho}, p) \to \exists q(q < p \amp \omega_2(\gn{\rho}, q)))$.
Поэтому $W_2(\gn{\rho}, q)$ ложно при любом $q$. Значит, по
выразимости $W_2$, $\proves \neg\omega_2(\gn{\rho}, q)$ для любого $q$, в том
числе и для $q < t$. Отсюда можно показать, что $\proves \neg\omega_2(\gn{\rho},
0) \amp \neg\omega_2(\gn{\rho}, 1) \amp \dots \amp \neg\omega_2(\gn{\rho}, t-1)$.
Применив лемму мы получаем, что $\forall q(q < t \to \neg\omega_2(\gn{\rho},
q))$, от этого же можно перейти к $\neg\exists q(q < t \amp
\omega_2(\gn{\rho}, q))$.

Обратно, пусть $\proves \neg\rho(\gn{\rho})$. Пусть $t$~--- гёделев номер
доказательства. Раз так, то $W_2(\gn{\rho}, t)$ истинно. По непротиворечивости
формальной арифметики это значит, что $W_1(\gn{\rho}, p)$ при любом $p$ ложно
(иначе окажется, что найдутся как доказательство, так и опровержение
$\rho(\gn{\rho})$). Значит, доказуемо $\neg\omega_1(\gn{\rho}, p)$ при любом $p$
(т.е. никакой из $p$ не является доказательством $\rho(\gn{\rho})$). Как
частный случай, $\neg\omega_1(\gn{\rho}, x)$ доказуемо для всех $x$, не
превышающих $t$, поэтому $\proves \neg\omega_1(\gn{\rho}, 0) \amp
\neg\omega_1(\gn{\rho}, 1) \amp \dots \amp \neg\omega_1(\gn{\rho}, t)$. Отсюда
можно показать $\proves p \le t \to \neg\omega_1(\gn{\rho}, p)$.

Рассмотрим формулу $(p > t) \to \exists q(q < p \amp \omega_2(\gn{\rho}, q))$.
Формула утверждает следующее: «если некоторый $p$ больше $t$, то найдется такой
$q$, меньший $p$, что $W_2(\gn{\rho}, q)$». Очевидно, что данная формула
истинна, ведь если мы возьмем $t$ в качестве такого $q$, то $W_2(\gn{\rho}, t)$
истинно по предположению. В силу выразимости $W_2$ в формальной арифметике
формула также и доказуема.

Легко показать, что из этих утверждений и из того, что $p \le t \vee p > t$,
можно вывести $\neg\omega_1(\gn{\rho}, p) \vee \exists q(q < p \amp
\omega_2(\gn{\rho}, q))$, а отсюда~--- $\forall p(\omega_1(\gn{\rho}, p) \to
\exists q(q < p \amp \omega_2(\gn{\rho}, q)))$, то есть $\proves \rho(\gn{\rho})$.
Однако, мы предположили $\proves \neg\rho(\gn{\rho})$, и исходя из него,
вывели $\rho(\gn{\rho})$, т.е. показали противоречивость формальной арифметики.
Значит, $\neg\rho(\gn{\rho})$ также недоказуемо, если арифметика
непротиворечива.
\end{proof}

Выберем утверждение, которое показывает непротиворечивость арифметики, т.е.
показывает отсутствие такой формулы $\varphi$, что и $\varphi$ и $\neg\varphi$
доказуемы. Поскольку в противоречивой теории можно вывести любое утверждение,
нам достаточно проверить это для какого-то конкретного $\varphi$, пусть это
будет $1 = 0$. Ясно, что $\neg 1 = 0$ доказуемо. Тогда для доказательства
непротиворечивости арифметики нам достаточно доказать $\forall
p(\neg\mathrm{Proof}(\gn{1 = 0}, p))$. Обозначим это утверждение за
$\mathrm{Consis}$.

\begin{thm}[Вторая теорема Гёделя о неполноте арифметики]
Если арифметика непротиворечива, то в ней не существует доказательства
$\mathrm{Consis}$.
\end{thm}
\begin{proof}
Рассмотрим формулу $\mathrm{Consis} \to \sigma(\gn{\sigma})$. Данная формула в
точности соответствует условию первой части первой теоремы Гёделя о неполноте
арифметики: если арифметика непротиворечива, то самоприменение формулы $\sigma$
истинно, т.е.\ недоказуемо; напомним, что $\sigma(x) \equiv \forall
p\neg\omega_1(x, p)$.

Рассуждение, доказывающее теорему Гёделя, можно формализовать, получив
доказательство данной импликации. Теперь, если у нас будет доказательство
утверждения $\mathrm{Consis}$, то по правилу Modus Ponens мы также получаем
доказательство утверждения $\sigma(\gn{\sigma})$, что невозможно по первой
теореме Гёделя.
\end{proof}

В данном месте есть очень поучительный пример, показывающий важность
формализации, и существенность деталей, которые мы опустили в доказательстве
второй теоремы о неполноте. В самом деле, вместо формулы $\mathrm{Consis}$ мы
могли бы попробовать рассмотреть формулу $\mathrm{Proof}_1(a) :=
\mathrm{Proof}(a, x) \amp \neg(\mathrm{Proof}(\gn{1 = 0}, x))$ и построить по ней
формулу $\mathrm{Consis}_1 := \forall x\neg\mathrm{Proof}_1(\gn{1 = 0}, x)$. На
первый взгляд, замена равноценна: действительно, если арифметика непротиворечива,
то тогда $\mathrm{Proof}_1$ ничем не отличается от $\mathrm{Proof}$, поскольку
$1 = 0$ тогда действительно недоказуема. Если же арифметика противоречива, то
$\mathrm{Consis}_1$ доказуема, как и любая другая формула.

Однако, $\mathrm{Consis}_1$ легко доказать. Пусть $\pi(x) \equiv
\mathrm{Proof}(\gn{0 = 1}, x)$, тогда:
\begin{align*}
&(1\dots n)\ \neg(\pi(x) \amp \neg(\pi(x))) &&\text{Доказуемо в и.в.}\\
&(n+1)\ \neg a = 0 &&\text{Аксиома А.x}\\
&(n+2)\ \neg(\pi(x) \amp \neg(\pi(x))) \to \neg a = 0 \to \neg(\pi(x) \amp \neg(\pi(x))) &&\text{Сх. акс. 1}\\
&(n+3)\ \neg a = 0 \to \neg(\pi(x) \amp \neg(\pi(x))) &&\text{M.P. } n, n+2\\
&(n+4)\ \neg a = 0 \to \forall x\neg(\pi(x) \amp \neg(\pi(x))) &&\text{Введение } \forall \text{ к } n+3\\
&(n+5)\ \forall x\neg(\pi(x) \amp \neg\pi(x)) &&\text{M.P. } n+1, n+4
\end{align*}
Получается, что мы здесь, продолжив рассуждение в соответствии со второй
теоремой, можем сразу получить противоречие и показать противоречивость
арифметики? На самом деле нет, поскольку мы некритично обобщили вторую теорему
на случай формулы $\mathrm{Consis}_1$, тогда как она существенно использует
внутреннее устройство $\mathrm{Consis}$. У $\mathrm{Consis}_1$ ведь есть
«слепое пятно»~--- формула $1 = 0$, на которой его результат не вычисляется, а
постулируется, и это значительное отличие.

Чтобы абстрагироваться от конкретного вида $\mathrm{Consis}$, Гильбертом и
Бернайсом были предложены следующие условия выводимости, позволяющие считать,
что формула действительно выражает непротиворечивость арифметики. Мы приведем их
в более позднем варианте, сформулированном Лёфом.

Пусть $\mathrm{Consis} \equiv \neg\mathrm{Provable}(\gn{1 = 0})$. Тогда
$\mathrm{Provable}$ должно отвечать следующим свойствам:
\begin{enumerate}
\item $\proves \alpha$ влечет $\proves \mathrm{Provable}(\gn{\alpha})$
\item $\proves \mathrm{Provable}(\gn{\alpha}) \to
  \mathrm{Provable}(\gn{\mathrm{Provable}(\gn{\alpha})})$
\item $\proves \mathrm{Provable}(\gn{\alpha \to \beta}) \to
  \mathrm{Provable}(\gn{\alpha}) \to \mathrm{Provable}(\gn{\beta})$
\end{enumerate}

Заметим, что в случае $\mathrm{Provable}_1(a) = \exists x(\mathrm{Proof}(a, x)
\amp \neg\mathrm{Proof}(\gn{1 = 0}, x))$ у нас будут сложности с демонстрацией
третьего правила. Анализ таблицы истинности импликации (вместе с полнотой и.в.)
указывает, что чтобы показать
\[ \proves \mathrm{Provable}_1(\gn{2 = 0 \to 1 = 0}) \to \mathrm{Provable}_1(\gn{2 = 0}) \to \mathrm{Provable}_1(\gn{1 = 0}) \]
нам надо показать $\proves \neg\mathrm{Provable}_1(\gn{2 = 0})$, поскольку
$\proves \mathrm{Provable}_1(\gn{2 = 0 \to 1 = 0})$ и $\proves
\neg\mathrm{Provable}_1(\gn{1 = 0})$. То есть, мы получим эту импликацию, только
если покажем доказуемость непротиворечивости формальной арифметики (ведь
$\proves \neg\mathrm{Provable}_1(\gn{2 = 0})$~--- это тоже вариант
$\mathrm{Consis}$).

Если бы речь шла об обычной формуле $\mathrm{Provable}$, мы бы решили вопрос
следующим неформальным, но формализуемым доказательством: если существует $x$,
такой, что $\mathrm{Proof}(\gn{2 = 0}, x)$, и существует $y$, что
$\mathrm{Proof}(\gn{2 = 0 \to 1 = 0}, y)$, то выполнено и $\mathrm{Proof}(\gn{1 =
0}, x \mathbin{@} y \mathbin{@} \gn{1 = 0})$.

В случае же с $\mathrm{Provable}_1$ нам все равно необходимо показывать $\proves
\neg\mathrm{Provable}_1(\gn{1 = 0})$.

Мы можем, конечно, указать общее соображение: если бы было верно
$\mathrm{Provable}_1(\gn{2 = 0})$, то тогда было бы верно $\proves 2 = 0$,
значит, теория противоречива и найдется доказательство чего угодно, в том числе
и $\proves \mathrm{Provable}_1(\gn{1 = 0})$, значит, выполнено
$\neg\mathrm{Provable}_1(\gn{2 = 0})$. Но это рассуждение касается только
значений формулы $\mathrm{Provable}_1$, сформулировано на мета-языке, и
непонятно, как из него сделать формальное доказательство $\proves
\neg\mathrm{Provable}_1(\gn{2 = 0})$.

\section{Теория множеств}
Теория множеств~--- это еще одна теория первого порядка, с которой мы
познакомимся в этом курсе. Мы добавим к исчислению предикатов один новый
двуместный предикат~--- отношение принадлежности $\in$.

Для изучения теории множеств мы введем несколько сокращений записи. Первое
сокращение~--- это новая логическая связка эквивалентность.

\begin{dfn}[Эквивалентность]
Запись $a \leftrightarrow b$ является сокращением записи для $a \to b \amp b \to
a$.
\end{dfn}

\begin{dfn}
Будем говорить, что множество $x$ является подмножеством множества $y$, если
любой элемент $x$ принадлежит $y$. Формально: $x \subseteq y$ является
сокращением записи для $\forall z(z \in x \to z \in y)$.
\end{dfn}

\begin{dfn}[Принцип объемности]
Два множества называются равными, если они являются подмножествами друг друга.
Формально: $x = y$ является сокращением записи для $x \subseteq y \amp y
\subseteq x$.
\end{dfn}

\begin{axm}[Аксиома равенства]
Равные множества содержатся в одних и тех же множествах. $\forall x\forall
y\forall z(x = y \amp x \in z \to y \in z)$.
\end{axm}

\begin{axm}[Аксиома пары]
Каковы бы ни были два различных множества $x$ и $y$, существует множество,
состоящее в точности из них. Будем записывать его так: $\{x, y\}$. Формально:
$\forall x\forall y(\neg x = y \to \exists p(x \in p \amp y \in p \amp \forall
z(z \in p \to (z = x \vee z = y))))$.
\end{axm}

\begin{axm}[Аксиома объединения]
Для любого множества $x$, содержащего хотя бы один элемент, найдется такое
множество, которое состоит в точности из тех элементов, из которых состоят
элементы $x$. Будем записывать его так: $\bigcup x$. Формально: $\forall
x(\exists y\, y \in x \to \exists p\forall y(y \in p \leftrightarrow \exists
s(y \in s \amp s \in x)))$.
\end{axm}

\begin{axm}[Аксиома степени]
Каково бы ни было множество $x$, существует множество $2^x$, содержащее в
точности все возможные подмножества множества $x$. Формально: $\forall
x\exists p\forall y(y \in p \leftrightarrow y \subseteq x)$.
\end{axm}

\begin{axm}[Схема аксиом выделения]
Для любого множества $x$ и любой формулы от одного аргумента $\varphi(y)$,
такой, что $b$ в нее не входит свободно, найдется такое множество $b$, в
которое входят те и только те элементы из множества $x$, что $\varphi(y)$
истинно. Формально: $\forall x\exists b\forall y(y \in b \leftrightarrow (y \in
x \amp \varphi(y)))$
\end{axm}

\begin{dfn}
Пересечением множеств $x$ и $y$ называется множество, состоящее в точности из
тех элементов, которые присутствуют и в $x$ и в $y$. Формально: $x \cap y$~---
это такое множество $z$, что $\forall t(t \in z \leftrightarrow t \in x \amp t
\in y)$
\end{dfn}

\begin{dfn}
Пустое множество $\varnothing$~--- множество, которому не принадлежит никакой
элемент: $\forall x\,\neg x \in \varnothing$.
\end{dfn}

\begin{thm}
\begin{enumerate}
\item Для любого множества $X$ существует множество $\{X\}$, содержащее в
  точности $X$.
\item Если существует хотя бы одно множество, то существует пустое множество.
\item Пустое множество единственно.
\item Для двух множеств существует множество, являющееся их пересечением.
\end{enumerate}
\end{thm}

\begin{dfn}
Дизъюнктным (разделённым) множеством называется множество, элементы которого не
пересекаются. Формально:
\[ \mathrm{Dj}(x) \equiv \forall y\forall z((y \in x \amp z \in x \amp \neg y = z) \to \neg\exists t(t \in y \amp t \in z)) \]
\end{dfn}

\begin{dfn}
Прямым произведением дизъюнктного множества $a$ называется множество $\times a$
всех таких множеств $b$, что:
\begin{itemize}
\item $b$ пересекается с каждым из элементов множества $a$ в точности в одном
  элементе
\item $b$ содержит элементы только из $\bigcup a$.
\end{itemize}
Формально:
\[ \forall b(b \in \times a \leftrightarrow (b \subseteq \bigcup a \amp \forall y(y \in a \to \exists!x(x \in y \amp x \in b)))) \]
\end{dfn}

Заметим, что прямое произведение дизъюнктного множества существует по аксиоме
выделения, поскольку $\times a \subseteq \bigcup a$.

\begin{axm}[Аксиома выбора]
Прямое произведение непустого дизъюнктного множества, не содержащего пустых
элементов, непусто. Формально:
\[ \forall t(\mathrm{Dj}(t) \to \forall x(x \in t \to \exists p(p \in x)) \to \exists p(p \in \times t)) \]
\end{axm}

\begin{axm}[Аксиома бесконечности]
Существует множество $N$, такое, что:
\[ \varnothing \in N \amp \forall x(x \in N \to x \cup \{x\} \in N) \]
\end{axm}

\begin{axm}[Аксиома фундирования]
В каждом непустом множестве найдется элемент, не пересекающийся с исходным
множеством.
\[ \forall x(x = \varnothing \vee \exists y(y \in x \amp y \cap x = \varnothing)) \]
\end{axm}

Аксиома фундирования исключает множества, которые могут принадлежать сами себе
(возможно, через цепочку принадлежностей): $X \in Y \in Z \in X$.

\begin{dfn}[Упорядоченная пара]
Упорядоченной парой двух множеств $a$ и $b$ назовем множество $\{a, \{a,
b\}\}$, еще будем записывать его так: $\langle a, b\rangle$.
\end{dfn}

\begin{lmm}
Упорядоченную пару можно построить для любых множеств, также $\langle a,
b\rangle = \langle c, d\rangle$ тогда и только тогда, когда $a = c$ и $b = d$.
\end{lmm}

\begin{axm}[Схема аксиом подстановки]
Если задана некоторая функция $f$, представимая в исчислении предикатов (то есть,
задана некоторая формула $\varphi$, такая, что $f(x) = y$ тогда и только тогда,
когда $\varphi(x, y) \amp \exists!z\,\varphi(x, z)$), то для любого множества $S$
существует множество $f(S)$~--- образ множества $S$ при отображении $f$.

Формально:
\[ \forall s(\forall x\forall y_1\forall y_2(x \in s \amp \varphi(x, y_1) \amp \varphi(x, y_2) \to y_1 = y_2) \to \exists t\forall y(y \in t \leftrightarrow \exists x(x \in s \amp \varphi(x, y)))) \]
\end{axm}

Через данную аксиому легко выразить аксиому выделения, однако аксиома подстановки
не может быть сведена к аксиоме выделения.

Данная аксиома позволяет строить множества большой и очень большой мощности:
допустим, если при $D(m) = 2^m$ положить $f(x) = D^x(\omega)$ (то есть $f(1) =
2^\omega$, $f(2) = 2^{2^\omega}$, $\dots$), то мощность $\bigcup f(\omega)$
будет больше мощности любого кардинального числа с конечным номером. Впрочем,
эта аксиома нужна даже для доказательства существования $2 \cdot \omega$:
аксиома бесконечности гарантирует существование только $\omega$.

\begin{dfn}[Бинарное отношение]
Бинарным отношением на множестве $X$ назовем подмножество множества всех
упорядоченных пар элементов из $X$.
\end{dfn}

На бинарных отношениях естественным образом вводятся отношения рефлексивности,
симметричности и транзитивности.

\begin{dfn}[Упорядочивание]
Отношение $R$ на множестве $S$ упорядочивает $X$, если это отношение
транзитивно и оно образует линейный порядок (строгое неравенство: справедливо
$\forall x\forall y(x \in X \to y \in X \to R(x, y) \vee x = y \vee R(y, x))$ и
$\forall x\,\neg R(x, x)$). Отношение вполне упорядочивает $S$, если к тому же
для любого непустого подмножества $B \subseteq S$ выполнено $\exists x(x \in B
\amp \forall y(y \in B \to \neg R(y, x)))$.
\end{dfn}

Также можно ввести понятие максимума, минимума, верхней грани, супремума.

\begin{dfn}
Множество $x$~--- транзитивное, если $z \in y,\ y \in x \to z \in x$.
\end{dfn}

\begin{dfn}
Ординал (порядковое число)~--- транзитивное, вполне упорядоченное с помощью
отношения $(\in)$ множество.
\end{dfn}

Ординалы~--- это некоторое обобщение понятия натурального числа. Для натуральных
чисел есть прямой аналог: $0 := \varnothing$; $1 := \{\varnothing\}$; $2 := 1
\cup \{1\}$ и т.п. Существование каждого из этих ординалов легко доказать.

\begin{dfn}
Ординал $x$ называется предельным, если $\neg x = \varnothing \amp
\neg\exists y(y \cup \{y\} = x)$.
\end{dfn}

\begin{dfn}
Ординал $x$ называется конечным, если он меньше любого предельного ординала. То
есть, он не содержит ни одного предельного: $\neg\exists t(t \in x \amp
\neg\exists y(y \cup \{y\} = t))$.
\end{dfn}

\begin{dfn}
Ординал $\omega$~--- это минимальный предельный ординал.
\end{dfn}

Ясно, что любое натуральное число меньше, чем $\omega$.

\begin{thm}
Ординал $\omega$ существует.
\end{thm}
\begin{proof}
Рассмотрим множество $N$, существование которого гарантирует аксиома
бесконечности. Заметим, что в этом множестве содержатся все конечные ординалы:
$\varnothing, \{\varnothing\}, \dots$ (это можно доказать по индукции на
мета-языке).

Для доказательства теоремы осталось выделить множество всех конечных ординалов
с помощью аксиомы выделения.
\end{proof}

Для ординалов можно определить арифметические операции.

\begin{dfn}
Верхней гранью подмножества $x$ некоторого вполне упорядоченного множества $S$
мы назовем всякий такой элемент $y$, что он больше всех элементов из $x$.

Минимальной верхней гранью подмножества $x$ некоторого вполне упорядоченного
множества $S$ мы назовем множество $\mathrm{Upb}_S(x) = \min\{y \mid y \in S
\amp \forall t(t \in x \to t < y)\}$.
\end{dfn}

Вообще, минимальная верхняя грань множества обычно носит название супремума
(обозначаемого как $\sup$), но в данном случае определение немного отличается:
мы указали строгое неравенство вместо нестрогого из-за чего минимальная верхняя
грань подмножества $\{1\}$ множества $\{1, 2\}$ не $1$ (как было бы естественно
для супремума), а $2$. Поэтому и обозначение (для избежания путаницы) также было
изменено: $\mathrm{Upb}$ вместо $\sup$.

Заметим, что $\mathrm{Upb}$ определена только для подмножеств некоторого
множества, (где минимум существует в силу вполне упорядоченности $S$, если для
$x$ в $S$ есть хотя бы одна верхняя грань). Однако, для ординалов, при условии
использования стандартного сравнения ординалов в качестве отношения $<$, мы
можем отказаться от такого ограничения.

\begin{dfn}[Арифметические операции над ординалами]
За $a + 1$ обозначим $x \cup \{x\}$. Тогда следующими рекурсивными определениями
мы введем операции сложения, умножения и возведения в степень:
\begin{align*}
a + b &\equiv
\begin{cases}
a, & b \equiv 0\\
(a + c) + 1, & b \equiv c + 1\\
\mathrm{Upb}_{\mathrm{ord}}\{a + c \mid c < b\}, & b \text{ --- предельный ординал}
\end{cases}\\
a \cdot b &\equiv
\begin{cases}
0, & b \equiv 0\\
(a \cdot c) + a, & b \equiv c + 1\\
\mathrm{Upb}_{\mathrm{ord}}\{a \cdot c \mid c < b\}, & b \text{ --- предельный ординал}
\end{cases}\\
a^b &\equiv
\begin{cases}
1, & b \equiv 0 \amp a > 0\\
(a^c) \cdot a, & b \equiv c + 1\\
\mathrm{Upb}_{\mathrm{ord}}\{a^c \mid c < b\}, & b \text{ --- предельный ординал}
\end{cases}
\end{align*}
\end{dfn}

Заметим, что в смысле введенной операции прибавления единицы, предельный
ординал~--- это такой ненулевой ординал $x$, для которого нет ни одного
ординала $y$, что $y + 1 = x$.

Так определенные операции на конечных ординалах ведут себя подобно обычным
арифметическим операциям на натуральных числах, однако в целом их поведение
может быть неочевидным. Например, легко заметить, что $1 + \omega = \omega$.

\begin{dfn}
Назовем множества $X$ и $Y$ равномощными, если найдется биективное отображение
$X$ на $Y$. Будем записывать это как $|X| = |Y|$. Будем говорить, что множество
$X$ имеет мощность не превышающую $Y$, если найдется инъективное отображение $X$
в $Y$. Будем записывать это как $|X| \le |Y|$. Будем записывать $|X| < |Y|$,
если известно, что $|X| \le |Y|$, но неверно, что $|X| = |Y|$.
\end{dfn}

\begin{dfn}[Кардинальные числа]
Кардинальное число~--- такой ординал $x$, что $y < x \leftrightarrow |y| <
|x|$.
\end{dfn}

Ясно, что все конечные ординальные числа являются кардинальными. Также,
например, $\omega$~--- кардинальное число (еще оно обозначается как $\aleph_0$,
если речь идет о мощности множеств). Кардинальное число $\aleph_1$
(соответствующая мощность называется континуум)~--- это следующее по порядку
кардинальное число за $\aleph_0$.

Верно ли, что $2^{\aleph_0} = \aleph_1$? Континуум-гипотеза (что это так) была
высказана довольно давно, и длительное время была одной из главных проблем в
теории множеств. Однако, оказалось, что данный факт не зависит от аксиоматики
Цермело--Френкеля, и можно как признавать этот факт, так и предполагать
существование промежуточных мощностей.

\section{Теорема Лёвенгейма--Сколема}
Как мы знаем из предыдущих разделов, рассмотрение формальной теории неполно,
если не сопровождается рассмотрением её моделей. Действительно, каковы могут
быть модели теории множеств? Например, обязан ли объект, соответствующий
множеству мощности $\aleph_1$ действительно иметь несчётное количество
элементов? В данном разделе мы прольём некоторый свет на этот вопрос.

\begin{dfn}[Мощность модели]
Пусть $M$~--- модель некоторой теории первого порядка. Напомним, что модель
задаётся предметным множеством $D$, и функциями, соответствующими всем
предикатам теории и всем функциональным символам теории. Тогда, назовем
мощность множества $D$ мощностью модели.
\end{dfn}

\begin{dfn}[Элементарная подмодель]
Пусть $M$~--- модель некоторой теории первого порядка, с предметным множеством
$D$, и пусть определено $D_1$, $D_1 \subset D$. Тогда структура $M_1$,
построенная на предметном множестве $D_1$ с предикатами и функциями,
получающимися из предикатов и функций $M$ путем сужения их области определения
на $D_1$ называется элементарной подмоделью $M$, если:
\begin{enumerate}
\item любая функция теории $f$ замкнута на $D_1$ (т.е.\ если $x_1 \in D_1,
  \dots x_n \in D_1$, то $f(x_1, \dots x_n) \in D_1$)
\item Любая формула $A(x_1, \dots x_n)$ теории при любых значениях $x_1, \dots
  x_n$ из $D_1$, истинная в $M$, истинна также и в $M_1$.
\end{enumerate}
\end{dfn}

Заметим, что второе условие необходимо: пусть некоторая теория 1-го порядка
содержит предикат равенства. Тогда формула $\exists x\exists y\,\neg(x = y)$ при
сужении предметного множества до одноэлементного перестаёт быть верной.

\begin{lmm}
Элементарная подмодель теории является моделью данной теории.
\end{lmm}
\begin{proof}
Рассмотрим некоторую доказуемую формулу теории $A$. Она является общезначимой в
$M$, значит, она является общезначимой в $M_1$.
\end{proof}

\begin{dfn}
Назовём теорию счётно-аксиоматизируемой (конечно-аксиоматизируемой), если ее
множество аксиом и правил вывода имеет счётную (конечную) мощность.
\end{dfn}

Заметим, что аксиомы и правила вывода явно содержат все «содержательные»
формулы в теории~--- формулы, которые могут использоваться в доказательствах хоть
в каком-нибудь качестве. В наших определениях ничего не запрещает создать
теорию, в которой заданы и дополнительные функции и предикаты, ни разу не
упоминаемые в аксиомах и правилах вывода, но таким предикатам и функциям мы
можем приписать произвольный смысл (например, все функции возвращают некоторую
заранее выбранную константу $c_0$, а все предикаты ложны), не повредив
корректности модели. Поэтому мы без ущерба для общности можем предположить, что
язык теории не содержит «несодержательных» формул. При этом заметим, что в силу
конечного размера каждой аксиомы и правила вывода любая конечно(счётно)-
аксиоматизируемая теория имеет конечное (счётное) множество всех возможных
«содержательных» формул.

\begin{thm}[Теорема Лёвенгейма--Сколема]
Пусть $M$~--- модель некоторой теории первого порядка, и пусть $T$~--- множество
всех формул этой теории. Тогда у $M$ есть элементарная подмодель, такая, что
$|M^*| = \max(|T|, \aleph_0)$.
\end{thm}
\begin{proof}
Для построения требуемой модели нам необходимо построить предметное множество
требуемой мощности и показать, что сужение теории на него даст нам элементарную
подмодель. Отсюда доказательство естественным образом разбивается на две части.

\emph{Часть 1. Построение множества.} Рассмотрим некоторое предметное множество
$D_0$. На его основе построим множество $D_0'$, рассмотрев все формулы из $T$, и
по каждой формуле (возможно) добавив некоторое количество элементов к $D_0$.

Фиксируем некоторую $n$-местную формулу $A(y, x_1, \dots x_n)$ из $T$.
Фиксируем некоторый набор аргументов $x_1, \dots x_n$ из $D_0$. В ходе
вычисления $A$ в результате применения функций к аргументам мы можем выйти из
$D_0$~--- добавим эти новые константы к $D_0'$. Их будет не более чем конечное
количество. Заметим, что формула $A$ может быть:
\begin{itemize}
\item тождественно истинной или ложной при любом $y$ из $D$~--- пропустим эту
  формулу и пойдем дальше. Поступим так же, если функция может быть как истинной
  так и ложной при разных значениях $y$ из $D_0$.
\item Найдутся такие $y$, что $A(y, x_1, \dots x_n)$ истинна, но при этом для
  любого $y$ из $D_0$ формула $A(y, x_1, \dots x_n)$ ложна~--- тогда добавим
  какой-нибудь один из этих $y$ к множеству $D_0'$. Также пополним множество
  всеми константами, которые необходимо добавить вследствие вычисления $A$ на
  данных аргументах. Этим мы также увеличим $D_0'$ на конечное количество
  элементов.
\end{itemize}

Рассмотрим множество $D_0$, $D_0 \subseteq D$, такое, что в него входят только те
элементы предметного множества, которые соответствуют константам (т.е.
нуль-местным функциям), упоминающимся в формулах из $T$. Если это множество
получается пустым~--- добавим к нему какую-нибудь одну константу из $D$. Оно
ляжет в начало счётной последовательности из множеств, $D_0 \subseteq D_1
\subseteq D_2 \subseteq \dots$ (такой, что $D_k' = D_{k+1}$), объединив которую,
мы получим множество $D^*$, которое и будет требуемым предметным множеством.

Заметим, что в результате однократного процесса пополнения мы увеличим
множество $D$ не более чем на $|T| \cdot \aleph_0 \cdot |D_0'|$ элементов (по
каждой формуле по каждой константе из $D_0'$ добавим не более, чем конечное
количество элементов). Поэтому мы не выйдем из мощности $|T|$, если $T$
несчётен, или из $\aleph_0$, если $T$~--- конечно. Значит, мощность $D^*$
также будет отвечать условию теоремы.

\emph{Часть 2.} Покажем, что структура $M^*$, полученная сужением модели $M$ на
предметное множество $D^*$, является элементарной подмоделью. Докажем это
индукцией по структуре формул. Рассмотрим некоторую формулу $A(x_1, \dots
x_n)$, при этом $x_i \in D^*$.

\emph{База.} Эта формула~--- предикат $P(f_1(x_1, \dots x_n), \dots f_k(x_1,
\dots x_n))$. Если $x_1, \dots x_n$ взяты из $D^*$, то все они добавлены на
некотором шаге, значит, найдётся такой $t$, что $x_i \in D_t$. Значит, все
результаты функций $f_j(x_1, \dots x_n)$ лежат в $D_{t+1}$. Значит, формула будет
определена на этих значениях. Очевидно, что истинность (ложность) формулы
сохранится.

\emph{Переход.} Пусть есть некоторая формула $A$~--- связка одной или двух
подформул, и $A$ имеет следующий вид:
\begin{enumerate}
\item $X \amp Y$, $X \vee Y$, $\neg X$, $X \to Y$~--- ясно, что все эти правила
  работают на сужении исходной модели, и истинность формулы сохраняется.
\item $\exists y\,B(y, x_1, \dots x_n)$. Фиксируем некоторый набор $x_1, \dots
  x_n$ из $D^*$. Пусть формула $A$ истинна в $M$, покажем её истинность в
  $M^*$. Заметим, что каждый из $x_i$ был добавлен в $D^*$ на конкретном шаге,
  то есть найдется такой номер $t$, что все $x_i$ принадлежат $D_t$. Тогда, по
  построению множества $D_{t+1}$, в нём неизбежно найдётся такой $y$, что $B(y,
  x_1, \dots x_n)$ определено и выполнено в $M$. Значит, $B$ будет выполнено и в
  $M^*$ (по предположению индукции). Значит, $A$ истинна и в $M^*$. Обратное
  также очевидно: если $y$ найдётся в узкой модели, то он тем более найдётся в
  широкой.
\item $\forall y\,B(y, x_1, \dots x_n)$. Доказательство зеркально соответствует
  предыдущему пункту. Истинность $\forall y\,B(y, x_1, \dots x_n)$ в $M$
  очевидно влечет истинность $\forall y\,B(y, x_1, \dots x_n)$ в $M^*$, так как
  $D^* \subseteq D$. Если же $\forall y\,B(y, x_1, \dots x_n)$ ложно в $M$, то
  существует такой $y$ из $D$, что $\neg B(y, x_1, \dots x_n)$ истинно в $M$,
  значит, по построению $D^*$ и в нём найдется элемент $z$ с таким свойством, то
  есть $B(z, x_1, \dots x_n)$ ложно в $M$. И по предположению индукции $B(z,
  x_1, \dots x_n)$ ложно в $M^*$. То есть $\forall y\,B(y, x_1, \dots x_n)$
  ложно в $M^*$.
\end{enumerate}
\end{proof}

Данная теорема в частности показывает, что для теории множеств (очевидно,
счётно-аксиоматизируемой) имеется счётная модель. В частности, это приводит к
следующему утверждению, называемому парадоксом Сколема: для любого множества,
определимого в аксиоматике Цермело--Френкеля, мы можем предложить способ
перенумеровать его элементы натуральными числами. В том числе это справедливо и
для несчётных множеств~--- то есть для тех множеств, для которых доказано
отсутствие этого способа.

Парадокс, тем не менее, здесь кажущийся, поскольку тут (некорректно) сталкивается
утверждение, формулируемое на предметном языке (об отсутствии способа пересчёта
внутри теории), с утверждением на мета-языке (о существовании этого способа вне
теории).

\begin{thebibliography}{8}
\bibitem{kleene1} С. Клини. Математическая логика --- М.: Изд-во «Мир», 1973
\bibitem{plisko} В. Е. Плиско, В. Х. Хаханян. Интуиционистская логика --- М.:
  Изд-во при мех.-мат. ф-те МГУ, 2009. Также доступно по ссылке
  \texttt{http://lpcs.math.msu.su/\string~plisko/intlog.pdf}
\bibitem{shen2} Н. К. Верещагин, А. Шень, Лекции по математической логике и
  теории алгоритмов, Языки и исчисления --- МЦНМО, 2002. Также доступно по
  ссылке \texttt{http://www.mccme.ru/free-books/shen/shen-logic-part2.pdf}
\bibitem{shen3} Н. К. Верещагин, А. Шень, Лекции по математической логике и
  теории алгоритмов, Вычислимые функции --- МЦНМО, 2002. Также доступно по
  ссылке \texttt{http://www.mccme.ru/free-books/shen/shen-logic-part3.pdf}
\bibitem{mendelson} Э. Мендельсон. Введение в математическую логику --- М.:
  Изд-во «Наука», 1971.
\bibitem{kleene2} С. Клини. Введение в метаматематику --- М.: Изд-во
  «Иностранная литература», 1957.
\bibitem{fraenkel} А. А. Френкель, И. Бар-Хиллел. Основания теории множеств ---
  М.: Изд-во «Мир», 1966.
\bibitem{cohen} П. Дж. Коэн. Теория множеств и континуум-гипотеза --- М.: Изд-во
  «Мир», 1969.
\end{thebibliography}

\end{document}
